%Subject: Paper for Special Issue on
%         Validated numerical methods and computer algebra


\documentstyle{jsc}

\title[Normal Form Method and Solutions of Nonlinear ODEs]
{Normal Form Method and Seminumerical Solutions of Nonlinear ODEs
                 \thanks{Supported by CNPq Brazil, project \#300894/93-7.}}

\author[V.F.Edneral]{Victor F. Edneral%\thanks{On leave of absence from
\\
%Grupo de Computa\c c\~ao Alg\'ebrica e Algoritmos \\
%        Instituto de Matematica - UFRGS \\
%        91501-970 Porto Alegre, RS, Brazil\\
%        e-mail: edneral@mat.ufrgs.br}
         Institute for Nuclear Physics\\
         of Moscow University,\\
         Leninskie Gory,\\
         Moscow, 119899, Russia\\
         e-mail: edneral@thp.npi.msu.su}

\date{}

\begin{document}

\maketitle


\section{Introduction}

\noindent
    A normal form method is based on a transformation of ODEs
system to a simpler set called the normal form. The importance of
this method for an investigation of ODEs near a stationary point
has been recognized for a long time. For the history of this subject
see, for instance, (Arnold, Anosov, 1988) or (Guckenheimer, Holmes, 
1986). Definitions of a normal form and a normalizing transformation
can be formulated in different approaches  for general and for
special cases. So there are very developed approaches for Hamiltonian 
systems, see, for example, (Deprit, 1969), (Hori, 1966), (Bruno, 1972), 
(Bruno, 1990, chapters 1,2). There are also many algorithms (and their 
realizations) for creating of normal forms and corresponding 
transformations. For the Hamiltonian case see an improved algorithm of 
Deprit and Hori in (Mersman, 1970) and its symbolic algebra realization
under REDUCE system in (Shevchenko, Sokolsky, 1993), a matter of a 
numerical creation of normal forms of Hamiltonians is described in 
(Godziewski, Maciejewski, 1990). Questions of integrability and 
convergence of the normalizing transformation in Hamiltonian case are
discussed in (Bruno, 1972, 1990), (Ito, 1989, 1992). Concerning 
algorithms for the creation of a normal form in a general case we 
mention here (in an addition to Bruno's books) papers (Walcher 91,93, 
Vallier, 1993). A connection between symmetry properties and convergence 
of a normalizing transformation have been investigated in (Bruno, 
Walcher, to be printed).

   In this paper we will use the algorithm based on the approach 
which was developed by A.D.Bruno (Bruno, 1971, 1972, 1989) for 
Poincare--Dulac's normal form. The 
important advantage of this approach is a possibility to 
investigate a wide class of equations in the united, easy 
algorithmized frame. In particular it provides a constructive way 
for a creation of approximations of families of 
periodic and conditionally periodic solutions in a form of 
power/Fourier series for real families or in a form of power
series in time dependent exponents for complex families. For this 
paper it is especially important that a problem of analyticity of 
used transformations is investigated and we are guaranteed 
that a domain of a convergency of corresponding (infinite 
as a rule) series is a nonzero one. The latest circumstance allows 
us to calculate 
periodic solutions and corresponding frequencies near stationary 
points with acceptable precision by finite formulaes and to be sure 
that with guaranty we can increase this precision by a calculation 
of highest terms. Except solutions themselves we can get origin 
conditions, which initiate periodic solutions in the form
of series in parameters or/and (for Hamiltonian systems) in the 
terms of the mechanical energy. I.e. we can produce some elements 
of a phase analyses of the systems.

   Another advantage of the using approach is an algorithmic 
simplicity of a creation of the normal form and the corresponding
transformation. We have for this procedure a direct recurrent 
formulae. Its usage does not demand keeping of some large 
intermediate results as it is in other algorithms. The approach is 
free from a necessity to solve intermediate systems of 
linear equations and from restrictions on low resonance cases.

   Below we describe a creation of the normal form and its 
application to a creation of periodic solutions of well known
equations. We assist our consideration with concrete examples 
produced by the NORT program  (Edneral, Khrustalev, 1985, 1992)
or by NORT and REDUCE system. The intention of the author is a
demonstration of the method on so transparent as it is possible
examples. Specularity of this paper is also in an attempt to
use a common approach to Hamiltonian and non Hamiltonian systems.


\section{Applications of normal forms}

Let us consider a system of ordinary autonomic nonlinear
differential equations:

\begin{equation}
\label{eq-planar}
\frac{d \overline{X}}{dt} = \overline{\Phi}(\overline{X},
\overline{a})
\end{equation}

\noindent
here $\overline{\Phi} = \{\Phi_1, \ldots, \Phi_n\}$ - a vector of right 
sides of the system which are functions in $\overline{X} = \{X_1, 
\ldots, X_n\}$ --- variables of the system and in $\overline{a} = 
\{a_1, \ldots, a_m\}$ --- its parameters. It means that $\overline{X}$  
are unknown functions in $\overline{a}$ and time --- $t$.

Such a type of equations originates from many scientific and engineering
problems where oscillations, vibrations or wave processes take place.
The standard way of an investigation of such systems is:

\begin{enumerate}

\item
a bifurcation analysis, i.e. an investigation of a picture of a behavior
of system 
solutions in dependence on parameters. It is especially important
if sharp change of this behavior occurs at some their values;

\item
a phase picture, i.e. an investigation of a behavior of system 
solutions in dependence on initial conditions;

\item
a calculation of system solutions. 

\end{enumerate}

If we have analytic solutions of a system, we have a clear picture of
behavior of the system, but we can have such solutions very rarely.
As a rule we have numerical solutions, but a numerical investigation 
of the above items 1 and 2  is sometimes a complicate problem. It is not 
simple also to create numerical solutions in unstable cases. As a 
sequence we may be interested in a some "intermediate" approach which 
would be between analytic and numerical methods.

The normal form method widely used for a bifurcation analysis. About 
methods of such an investigation see, for example, (Hassard, Kazarinoff, 
Wan, 1981) and (Guckenheimer, Holmes, 1986).
At these books you can see that a numerical bifurcation analyses is 
based really on the normal form method. The main idea here is in
replacing of system (\ref{eq-planar}) by some "model" system with
polynomial right sides of a finite order and in recasting then low 
orders of them to the canonical (normal) form. From the lowest 
coefficients of the normal form we can make conclusions 
about a qualitative behavior of the origin system. It is enough to
know only lowest orders of the normal form for such an analysis. 
Sometimes this job is produced by hands, sometimes by computer algebra 
systems, see, for example, (Rand, Armbruster, 1987).

   Here we try to demonstrate, that calculation of high orders of
the normal form can be useful also for targets 2 and 3 mentioned above.
We restrict our consideration here mainly by a creation of approximations 
to periodic solutions but we have a positive material of using
this approach in a case of conditionally periodic solution.

\section{The normal form method}

   At the investigation of equations of type (\ref{eq-planar}) there
three preliminary steps are usually produced. The first step is a shift 
of variables of the system $\overline{X}$ in such a way that we 
have in new variables: $\overline{\Phi}(0,\overline{a}) = 0$. 
The constant zero vector will be a solution of the 
system. A corresponding point (zero here) is called an "immovable", 
"stationary" or "equilibrium" point. A system can have more then one 
such a point. Each of them is to be investigated separately.
The second step is a reduction of the system to the "model" 
form which approximates vector $\overline{\Phi}(\overline{X},
\overline{a})$ by a set of polynomials of a finite order. If 
$\overline{\Phi}$ is an analytical function of $\overline{X}$
in some domain of the stationary point it can be realised with any
desirable precision. Often this step is made simultaneously with a
reduction of the system to the central manifold. After that
right sides of the system will be polynomials without 
constant (absolute) terms. By a linear complex change of its variables it is 
possible to reduce a linear parts of the system to the canonical
form. After this (the third) step we have: 

\begin{equation}
\label{eq-prep}
\frac{d\overline{Y}_i}{dt}\ =\ \hat{A} \overline{Y}+
      \overline{F}(\overline{Y},
\overline{a})
\end{equation}

\noindent
where $\hat{A}$ matrix has a canonical (Jordan's) form and 
$\overline{Y}$ -- new variables of the system after the above 
transformations.

   In this paper we suppose also that matrix $\hat{A}$ has a 
diagonal form. Really the described method is formulated without 
such a restriction. It has rather a technical character and can be 
eliminated with a development of the NORT program. We suppose also
that the dependence of the right side in parameters has a polynomial
character also. It has a more principal character. 

Finally we will have the system which is considered in this
paper (below for a simplification of a notation we will hint the 
dependence on parameters):

\begin{equation}
\label{eq-general}
\frac{dY_i}{dt}\ =\ \lambda_iY_i+F_i(Y_1,Y_2,\ldots,Y_n),
\quad \ i=1,\ldots ,n
\end{equation}

\noindent
here  $\{\lambda_i\}$ is a non-zero vector of eigenvalues of the linear 
parts of the right sides.  $\{F_i\}$ is a vector of polynomials of 
finite degrees without constant and linear terms. In other words we 
impose further restrictions:

\begin{description} 
\item[$\bullet$]
    we  take into account autonomic \footnote{Actually
periodic dependence on time can be eliminated by changing time
dependent functions for the proper solutions of additional ODEs. For 
instance $\cos(\omega t)$ can be simulated by a couple of additional 
autonomic equations with proper origin conditions.}
systems with a polynomial nonlinearity 

\item[$\bullet$]
    systems are  investigated near a stationary point,  i.e.  the 
vector $\{(0,\ldots,0)\}$ is a solution of the system;

\item[$\bullet$]
    we  suppose that a linear part of the right hand side can be 
    diagonalized  and this  operation  has  been  done (actually  
    the  algorithm  can  be  
generalized for a case of a triangular matrix (Bruno, 1988).
But anyway the eigenvalues must be know). We also suppose that 
its eigenvalues $\lambda_i$ have numerical values\footnote{This is 
not a very dramatic requirement, see, e.q., a treatment of van der 
Pol's equation below.}.

\end{description} 

    Let us now rewrite the equation (\ref{eq-general}) in the form:
    
\begin{equation}
\label{eq-bform}
\frac{dY_i}{dt}\ =\ \lambda_i Y_i+Y_i \sum_{\overline{Q}_i}
F_{i,\overline{Q}_i}\overline{Y}^{\overline{Q}_i},
\end{equation}

\noindent
where $\overline{Y}^{\overline{Q}_i}$ is defined as a product:

$$
\prod_{j=1}^{n} {Y_j}^{Q_{ij}}, \quad
Q_{ij}\ =\ \left\{ \begin{array}{ll}
                    0,1,2,\ldots & \mbox{if $i\not =j,$} \\
                    -1,0,1,\ldots & \mbox{if $i=j.$}
                   \end{array}
           \right.
$$

$Q_{ij}$ here are components of a vector $\overline{Q}_i=\{Q_{ij}\}$.
$Q_{ij}$ are non-negative integers  with the exception of 
$\{Q_{ii}\} \geq -1$,  because factor $Y_i$ has been carried out from 
sum in (\ref{eq-bform}). Summation in (\ref{eq-bform}) is over all 
possible $\overline{Q}_i$-vectors.

    In such terms a normalizing transformation may be written as:

\begin{equation}
\label{n-trans}
Y_i\ =\ Z_i + Z_i \ \sum_{\overline{Q}_i}
H_{i,\overline{Q}_i} \overline{Z}^{\overline{Q}_i}.
\end{equation}

    Coefficients $H$ in transformation (\ref{n-trans}) can be chosen so
that equation (\ref{eq-bform}) will have the form:

\begin{equation}
\label{eq-nform}
\frac{dZ_i}{dt} =  \lambda_i Z_i + Z_i \sum_{(\overline{Q}_i,
\overline{\lambda})=0} G_{i,\overline{Q}_i} \overline{Z}^{\overline{Q}_i}.
\end{equation}
\vspace{0.07in}

    The great  difference between (\ref{eq-nform}) and (\ref{eq-bform}) 
is a restriction on the domain of summation, which is defined by the 
equation:

\begin{equation}
\label{sc-prod}
(\overline{Q}_i,\overline{\lambda})\ \stackrel{\rm def}{=} \sum_{j=1}^{n} 
Q_{ij}\lambda_j\ =\ 0.
\end{equation}

   The form (\ref{eq-nform}) is called the normal form for 
(\ref{eq-bform}). The $H$ and $G$ coefficients in (\ref{n-trans})  
and (\ref{eq-nform}) may be found  from the recurrent formulae 
(Bruno, 1971, 1989):

\begin{equation}
\label{rec-norm-trans}
\begin{array}{ll}
&G_{i,\overline{Q}_i}  +  (\overline{Q}_i,\overline{\lambda})
H_{\mathstrut \mathstrut \mathstrut i,\overline{Q}_i}= \\
&-  \sum_{\mathstrut \mathstrut \mathstrut j=1}
        ^{\mathstrut \mathstrut \mathstrut n}
\sum_{\overline{P}+\overline{R}=\overline{Q}_i}
 (P_j+ \delta _{ij})H_{i,\overline{P}}G_{j,\overline{R}} \\
&+\Phi_{i,\overline{Q}_i}^{\mathstrut \mathstrut \mathstrut \ }.
\end{array}
\end{equation}

\noindent
Here the second summation in the right-hand side is over all vectors 
$\overline{P}$ and $\overline{R}$ for which  
$\overline{P}+\overline{R}=\overline{Q}_i$; $P_j$ in parentheses is the
$j$- component of vector $\overline{P}$; $\Phi_{i,\overline{Q}_i}$ is a
coefficient at the factor $Z_i \overline{Z}^{\overline{Q}_i}$ in the 
polynomial $F_i$ from (\ref{eq-bform}), arguments of which have been
transformed by 
(\ref{n-trans}) \footnote{Take into account that the scalar production 
$(\overline{Q}_i,\overline{\lambda})$ will appeare in denominators of 
$H$ coefficients with any possible values of $\overline{Q}$. If one of
eigenvalue has a symbolic value this scalar production 
will provides a series of (complex in general) poles in this 
symbolic value.}.

The uncertainty in (\ref{rec-norm-trans}) is fixed below by conventions:

\begin{equation}
\label{fix-trans}
\begin{array}{l}
\mbox{$H_{\mathstrut i,\overline{Q}_i}=0 
   \qquad$ if $\quad(\overline{Q}_i,\overline{\lambda})=0,$} \\
\mbox{$G_{i,\overline{Q}_i}^{\mathstrut \ }=0 
   \qquad$ if $\quad(\overline{Q}_i,\overline{\lambda})\not =0$}
\end{array}
\end{equation}


Sometimes the normalizing transformation fixed by (\ref{fix-trans})
is called a "basic transformation".


\subsection{the main algorithm}

The algorithm of the calculation of $G$ and $H$ is based on
(\ref{rec-norm-trans}),(\ref{fix-trans}). Let us
choose the representation of sets of coefficients $G_{i,\overline{Q}_i}$
and $H_{i,\overline{Q}_i}$ so that they would be collected in such 
groups where a group of order $n$ contains only elements with such 
vector-indexes $\overline{Q_i} = \{q_{i,j}\}$ that the sum of their 
components is equal to the order, i.e. $\sum_{j}q_{i,j}=n$ for
each $i$. You can calculate the group of the next order for
$G$ and $H$ by using groups of $G$ and $H$ with smaller order only, 
i.e. (\ref{rec-norm-trans}) is a recurrent formulae.

\vspace{0.07in}

The algorithm:

Let $n$ is a dimension of the system, for its normalization till
order m we are to do:

\begin{enumerate}
\renewcommand{\theenumi}{(\roman{enumi})}

\item
for $i=1,2,\ldots,n$ do:\\
Calculate all square in $Y_i$ elements in right sides of  
$F_i(\overline{Y})$, i.e. calculate the group of the first 
order of the set $\Phi_{i,\overline{Q}_i}$ and sort it into two 
sets in dependence on a value of 
scalar product (\ref{sc-prod}). The first set where this product 
is zero will be the first order group of $G_i$ and the second set
after a division on the value 
of the corresponding scalar product will be the first order group of $H_i$ 

\item
for $k=2,3,\ldots,m$ do:
   \begin{enumerate}

   \item
   for $i=1,2,\ldots,n$ do:\\
   calculate the group of order $k$ of the set
   $\Phi_i$ for which substitute the $Y_i$ evaluated by (\ref{n-trans}) 
   till order $k-1$ in the $F_i(\overline{Y}(\overline{Z}))$
   and choose coefficients at monomials 
   $Z_i \overline{Z}^{\overline{Q}_i}$ as $\Phi_{i,\overline{Q}_i}$;
   
   \item
   for $i=1,2,\ldots,n$ do:\\
   Calculate the groups of $G_i$ and $H_i$ of order $k$  by a division 
   the set $\Phi_i$ into two sets as in item (a). After that you will
   have the full group of order $k$ for $G_i$ and preliminary values
   of such a group for $H_i$, because those values of the vector 
   $\overline{Q}_i$ 
   which takes place at non-zero elements into the summation in
   (\ref{rec-norm-trans}) can not zero the scalar product,
   as a sum of a vector which zeroes this product ($\overline{R}$)
   and a vector which does not do it ($\overline{P}$).

   \item
   for $i=1,2,\ldots,n$ do:\\
   for $j=1,2,\ldots,n$ do:\\
   supplement the group of order $k$ of $H_i$ with properly sorted
   multiplications of {\em all} elements of such groups of $H_i$
   and $G_j$ so that their summary order (i.e. $\sum_{s=1}^n (P_s+R_s)$) 
   would be equal to $k$. These multiplications would be multiplied also 
   by a factor
   $(P_j+\delta_{i,j})$, which can be equal to zero and therefore
   should be calculated before the other factors. Before the supplement
   all elements are to be divided on the corresponding scalar products. 

   \end{enumerate}
\end{enumerate}

It is not so heavy to estimate a cost of this algorithm but here
such estimation has no a deep sense because of relative high
complexity of a calculation of right sides of the nonlinear system.
Under such circumstances it is very important to calculate these
parts very economic, using as possible the fact that we need at 
the each step of (ii) a single homogeneous 
"slice" of $\overline{F}$ of order $k$ only and its terms of lower orders 
will be not changing during the following operations. The problems 
with the optimization of the right sides evaluation is one of the main reason 
of an absence of closed interface for the NORT package. 

\subsection{periodic solutions, a convergence of series and adequacy
         of the normal form to the origin system}

  The sums in (\ref{n-trans})  and (\ref{eq-nform}) may 
include  an infinite  number  of terms for physically interesting 
cases even though the sum in (\ref{eq-bform}) has a finite number 
of ones. The question about properties of a convergence and a 
smoothness of transformation arises. For adequacy of the normal form 
to the origin system in a some domain of a stationary point we need 
to have analytical normalizing transformation. The problem of 
analyticity of (\ref{n-trans}) have been investigated in (Bruno, 1971, 
1989). The result is shortly that the transformation (\ref{n-trans})
is divergent in a general case, but there can be found some subsets of 
the formal invariant sets of the normalized system (\ref{eq-nform})
where the normalizing transformation is convergent and which are 
therefore adequate to some invariant subsets of the original system. 
Bruno calls them analytical sets. 

We will see below the subspace of a coordinates space of normalized 
system which includes only variables $Z_i$ for which $Re(\lambda_i) 
= 0$. Actually, if this is a null-space we have no deep problems 
with the convergence,
because in this case there is a result about topological equivalence 
of the original system to the system which is normalized
till a {\em finite} order \footnote{We will not discuss here this 
circle of questions and send a reader for the references to Bruno's 
books.}. The variables for which corresponding eigenvalues are not
pure imaginary increase to infinity (for $Re(\lambda_i)>0$) or
vanish (for $Re(\lambda_i)<0$) with time. Therefore we have for
periodic solutions only the above mentioned subspace. Let us extract
in this subspace the set ${\em A}$:

\begin{equation}
\label{cond-A}
\begin{array}{lll}
{\em A} = \{\overline{Z} : 
            &\psi_i = \lambda_i Z_i a, \ \ &if \ Re(\lambda_i) = 0;\\
            &Z_i = 0,                      &if \ Re(\lambda_i) \neq 0;
\quad i=1,\ldots,n \}
\end{array}
\end{equation}

\noindent
where $\psi_i$ is an element of a right side of the normalized system
and $a$ is a free parameter, the same for all values of $i$.

%\newtheorem{theorem}{Theorem}
%\begin{theorem}
{\bf Theorem 3.1}
   The set ${\em A}$ contains periodic and conditionally periodic 
solutions of the normalized system including stationary points.
%\end{theorem}

   Further let us collect all pure imaginary eigenvalues in groups with
the property: each group includes all $\lambda$'s which are co-measured
in pairs \footnote{I.e. a ratio of any non-zero elements within the
group is a rational number.}. Following Bruno we will speak about 
coordinate subspace which contains only the coordinates (variables) 
with numbers corresponding one of these groups as about the "rational"
subspace. Let us designate the intersection of the set ${\em A}$ with 
a union of all such rational subspaces by $\stackrel{\sim}{\em A}$.

%\newtheorem{theorem1}{Theorem}
%\begin{theorem1}
{\bf Theorem 3.2}
   The set $\stackrel{\sim}{A}$ are filled with periodic solutions 
of the normalized system including stationary points and there exists
an analytical invertible transformation which transfers the 
original system to the system which is normalized on the set
$\stackrel{\sim}{A}$.
%\end{theorem1}

The important sequence from the above theorems is that if we find
the set $\stackrel{\sim}{\em A}$ we find the periodical solutions of
the normalized system, which contracts to the stationary point. 
Along these solutions the normalizing transformation will be analytical. 
Thus the normalized system will adequate to the origin one along these 
solutions.

It is easy to see that a fulfillment of (\ref{cond-A}) is enough for 
integration of the system, because of $a$ will be a constant. It provides
also the way for building integrals of motion.

A set exists also which contains conditionally periodic solutions
such that the normalizing transformation is analytic along them.
The definition of such a set see in (Bruno, 1989).

\subsection{the general scheme of the ODEs investigation}

General scheme of an investigation of a nonlinear ODEs system 
by the normal form method near each stationary point is:

\begin{enumerate} 

\item
    recasting the system to a (model) polynomial form;

\item
    the linear normalization of the system, i.e. the reduction of 
a linear part of the right sides of the system to Jordan's or to a 
diagonal form and an investigation of the corresponding eigenvalues;

\item
    the non-linear normalization of the system,---the creation of
the normal form and the corresponding normalizing transformation;
(sometimes it is called recasting the system in a "resonance form").

\item
    a bifurcation analysis of the system in parameters by the 
investigation of lowest non-disappeared orders of the normal form;


\item
    building the periodic and conditionally periodic solutions
which contract to the stationary point;

\item
    decreasing an order of the normalized system if it necessary and
a repetition of the present investigation from item 2 near each 
stationary point of the system decreased.

\end{enumerate} 

\section{Second order systems}

All questions about convergence and integrability of normal forms for 
all second order systems of type \ref{eq-general} near a stationary 
point have been investigated by (Bruno, 1971,1989) for the case when 
the stationary point is an elementary especial point, i.e. both 
eigenvalues of the system are not equal zero simultaneously.
The result of this investigation shortly is: the normal forms for
equations defined as the above are integrable. For each case there are 
written integrals of the normalized system. The result regarding the
convergence is the following. 
Let $\lambda_2 \neq 0$ and $\lambda \stackrel{\rm def}{=}
 \lambda_1/\lambda_2$. If $Im(\lambda) \neq 0$ or if 
$\lambda > 0$ then the transformation is convergent and we can
produce an approximate solution of the origin system from known
integrals of the normal form by transformation (\ref{n-trans}) with 
any desirable precision. In other words cases of "focus" and "knot" 
can be treated without any additional demands. At a negative irrational
$\lambda$ the convergence will take place if for all non-zero vectors 
$\overline{Q}$ with integer elements there exist positive $\varepsilon,\nu$ 
such that $|(\overline{Q},\overline{\lambda})| > \varepsilon 
(|q_1|+|q_2|)^{-\nu}$. This condition can be checked before
creation of the normal form. The latest case is a one partial case
of a saddle point. 
But at real non-positive rational $\lambda = -m/n \leq 0$ we will have
convergence under some additional 
requirements on the normal form. It is interesting case. It includes
the cases of a center and a limit circle. Let us look at a couple
of examples.

\subsection{duffing's equation}

   This is an equation of the second order, which origins from a
problem of a mathematical pendulum:

\begin{equation}
\label{eq-pendulum}
\frac{d^{2}\phi}{dt^2}+\omega_0^2\sin(\phi) \ = \ 0
\end{equation}


\noindent 
where $\phi$ is an angle of deviation of the pendulum and 
which can be approximated by the form (units below have been chosen so 
that $\omega_0 = 1$):

\begin{equation}
\label{eq-pend}
\frac{d^{2}\phi}{dt^2} \ = \ - \phi + {1 \over 6}\phi^3
\end{equation}

By considering $\phi = \sqrt{6} X$ we get Duffing's equation:

\begin{equation}
\label{eq-duffing}
\frac{d^{2}X}{dt^2} \ = \ - X + X^3
\end{equation}

This is a Hamiltonian equation with the full energy $H$:

\begin{equation}
\label{H-duffing}
H = {1 \over 2}(\frac{dX}{dt})^2 + {1 \over 2} X^2 - {1 \over 4} X^4
\end{equation}


By linear complex changing for variables:

\begin{equation}
\label{line-norm}
\begin{array}{ccc}
X & = & Y_1 + Y_2 \\
\frac{\mathstrut dX}{dt}& = & i \ (Y_1 - Y_2)
\end{array}
\end{equation}

\noindent
we can rewrite it in the diagonalized form:

\begin{equation}
\label{planar-duff}
\begin{array}{lcrl}
\frac{\mathstrut dY_1}{\mathstrut dt}& = &  i &\!\! Y_1 - {i \over 2}
                                                       \ (Y_1 +Y_2)^3\\
\frac{\mathstrut dY_2}{\mathstrut dt}& = & -i &\!\! Y_2 + {i \over 2}
                                                       \ (Y_1 +Y_2)^3
\end{array}
\end{equation}


\noindent
Note that a couple of equations the above mentioned have complex conjugated 
coefficients at terms with exchange $Y_1 \leftrightarrow Y_2$ It will 
take place always if the origin system has real coefficients. 

   The vector of eigenvalues of 
a linear part of the system is $\overline{\lambda} = \{i,-i\}$.

In accordance with the definition of the normal form: 
$$
\frac{dZ_i}{dt} \ = \ \lambda_iZ_i + Z_i
                                   \!\!\!\!\!\!\!\!\!\!\!\!\!\!\!\!\!
                                   \!\!\!\!\!\!\!\!\!\!\!\!\!\!\!\!\!
                           \sum_{\begin{array}{c}
                                 p_i \geq -1 \\
                           p_1,\ldots,p_{i-1},p_{i+1},\ldots,p_n \geq 0 \\
                           (\overline{\lambda},\overline{p}) = 0 
                                              \end{array}}
                                   \!\!\!\!\!\!\!\!\!\!\!\!\!\!\!\!
                                   \!\!\!\!\!\!\!\!\!\!\!\!\!\!\!\!
                               g_{i,p_1,\ldots,p_n}
                               \prod_{k = 1}^{n} Z_k^{p_k},            
\quad i=1,\ldots ,n
$$

\noindent
we will have for the case $n = 2$ 
at sums in the right sides of equations only terms, where 
$(\overline{\lambda},\overline{p}) = i\ (p_1 - p_2) = 0$.
I.e. sums at the right sides of the normal form of Duffing's 
equation will include only terms where $p_1 = p_2$: 

\begin{equation}
\label{n-duffing}
\begin{array}{lcrl}
\frac{\mathstrut dZ_1}{\mathstrut dt}& = &  i &\!\! Z_1 
+ Z_1\ (g_{1,1,1} Z_1 Z_2 + g_{1,2,2} Z_{1}^2 Z_{2}^2 + \ldots)\\
\frac{\mathstrut dZ_2}{\mathstrut dt}& = & -i &\!\! Z_2
+ Z_2\ (g_{2,1,1} Z_1 Z_2 + g_{2,2,2} Z_{1}^2 Z_{2}^2 + \ldots)
\end{array}
\end{equation}

The condition (\ref{cond-A}) for the second order equation with 
eigenvalues $\lambda_1 = -\lambda_2$ has the form:


$$
\sum_{k=1,\ldots} g_{1,k,k} \ (Z_1 Z_2)^k = 
- \sum_{k=1,\ldots} g_{2,k,k} \ (Z_1 Z_2)^k
$$

As for any real (originally) equation of the second order the equalities:

\begin{equation}
\label{conj-2}
\begin{array}{c}
g_{1,i,i} = g_{2,i,i}^{*}\\
h_{1,i,j} = h_{2,j,i}^{*}
\end{array}
\end{equation}

takes place \footnote{The asterisk here notes a complex conjugation},
this condition for real equations of the second order has the form:

\begin{equation}
\label{cond-2}
\sum_{k=1,\ldots} Re(g_{1,k,k}) \ (Z_1 Z_2)^k = 0
\end{equation}

After a calculation of the normal form for Duffing's equation,
it can be seen that (\ref{cond-2}) is fulfilled identically
because all $g_{1,i,i} \mbox{ and } g_{2,i,i}$ are pure imaginary here. 
It means that in the case of Duffing's equation we have the
case which is usually  called the "center", when 
periodic solutions exist for any initial conditions.

Really, 
by multipling the first from equations (\ref{n-duffing}) by $Z_2$
and the second one by $Z_1$ we will have after their addition 
$\frac{d(Z_1 Z_2)}{dt} = 0$ and so (\ref{n-duffing}) has a trivial
form:

$$
\begin{array}{lccl}
\frac{\mathstrut dZ_1}{\mathstrut dt}& = &  &i \ \omega(Z_1 Z_2) \ Z_1 \\
\frac{\mathstrut dZ_2}{\mathstrut dt}& = & -&i \ \omega(Z_1 Z_2) \ Z_2  
\end{array}
$$

\noindent
because $\omega(Z_1 Z_2) \stackrel{\rm def}{=} 1 + g_{1,1,1}/i + 
g_{1,2,2}/i + \ldots$ 
is a real constant. Therefore the family of solutions of Duffing's 
equation in the normal form is:

$$ 
Z_1(t) = c_1 e^{+ i \omega(c_1 c_2) t}, \quad
Z_2(t) = c_2 e^{- i \omega(c_1 c_2) t}  
$$

\noindent
where $c_1, c_2$ are voluntary constants.

Now we can create the approximation to origin equation 
(\ref{eq-duffing}) by substitution the above found $Z_i$ into $Y_i$ 
with (\ref{n-trans})
and after that into $X$. If we choose complex conjugate values for 
$c_1, c_2$ we will have an approximation of the real solution in the 
form of truncated Fourier series.

   The calculation by the NORT package till $15^{\mbox th}$ order in $C = 
\sqrt{c_1 c_2}$ took 53 seconds on 33MHz 4Mbyte RAM  PC/AT-486DS 
computer in the symbolic mode of REDUCE 3.4.2. Due to equalities
(\ref{conj-2}) we do not need to calculate series for both of 
equations (\ref{planar-duff}). So we have 16 terms for  
each sums in normal form (\ref{n-duffing}) and 255 terms for
the each normalizing transformation from $Z_i$ to $Y_i$. The final
expression for a frequency has 16 terms and the expression for $X(t)$ 
has 141 terms. 

We got these series as series in the $C$ variable. It is not convenient
and for the final representation we calculated the $H$ as a series in 
$C$ by substituting found $X \mbox{ and } {dX \over dt}$ in
(\ref{H-duffing}). By inverting this series and substituting $C$ as 
a series of $H$ into expressions for the $\omega$ and $X$ we will
have the end result as series in $H$. For an economy of the place
we show here this result till the $5^{\mbox th}$ order only 
\footnote{If somebody
is interested in the full series in the REDUCE readable or FORTRAN
form, he can reach them via the name service by anonymous ftp from
euler.mat.ufrgs.br (143.54.2.220) in the directory pub/norm}:


\begin{equation}
\label{res-duffing}
\begin{array}{ll}

\omega =& 1
 -{\mathstrut 3 \over \mathstrut  4} H
 -{\mathstrut 69 \over \mathstrut  64} H^2
 -{\mathstrut 633 \over \mathstrut  256} H^3
 -{\mathstrut 110421 \over \mathstrut  16384} H^4
 -{\mathstrut 1318329 \over \mathstrut  65536} H^5 + \ldots\\

&\\
 
X =& \sqrt{2 H} \times \\
 &[ \ \ \cos(\omega t) (
   1
 +{\mathstrut 9 \over \mathstrut  16} H
 +{\mathstrut 271 \over \mathstrut  256} H^2
 +{\mathstrut 10779 \over \mathstrut  4096} H^3
 +{\mathstrut 243613 \over \mathstrut  32768} H^4
 +{\mathstrut 2963587 \over \mathstrut  131072} H^5)\\
 &-\cos(3 \omega t) H (
  {\mathstrut 1 \over \mathstrut  16}
 +{\mathstrut 3 \over \mathstrut  16} H
 +{\mathstrut 1209 \over \mathstrut  2048} H^2
 +{\mathstrut 127233 \over \mathstrut  65536} H^3
 +{\mathstrut 6907221 \over \mathstrut  1048576} H^4)\\
 &+\cos(5 \omega t) H^2 (
  {\mathstrut 1 \over \mathstrut  256}
 +{\mathstrut 11 \over \mathstrut  512} H
 +{\mathstrut 3107 \over \mathstrut  32768} H^2
 +{\mathstrut 25567 \over \mathstrut  65536} H^3)\\
 &-\cos(7 \omega t) H^3 (
  {\mathstrut 1 \over \mathstrut  4096}
 +{\mathstrut 1 \over \mathstrut  512} H
 +{\mathstrut 5805 \over \mathstrut  524288} H^2)\\
 &+\cos(9 \omega t) H^4 (
  {\mathstrut 1 \over \mathstrut  65536}
 +{\mathstrut 21 \over \mathstrut  131072} H)\\
 &-\cos(11 \omega t) H^5 (
  {\mathstrut 1 \over \mathstrut  1048576}) + \ldots \ ]

\end{array}
\end{equation}


The result of this calculation was verified  by two ways. The first 
one was a direct substitution of series (\ref{res-duffing}) in 
the original equation (\ref{eq-duffing}). After that we had only terms  
with negligible orders of $H$. The second way was a comparison of the
numerical solutions of Duffing's equation evaluated by Runge--Kutta 
method (by NAG's d02baf procedure with a 'tolerance' $= 10^{-12}$) and  
the values tabulated of the series the above mentioned at different 
values of $H$.
In view of (\ref{H-duffing}) we have $H=H(X=\phi/\sqrt{6},\frac{dX}{dt}
=\frac{d\phi}{dt}/\sqrt{6})$. On the other hand $H$ has its maximum 
physically reasonable value when at a zero velocity a maximum
deviation of a pendulum takes place, i.e. when $\phi = \pi/2$ and 
$\frac{d\phi}{dt}=0$. In this case we can choose $H \leq H_{max} 
\simeq 0.163$. Let us now introduce a function of a maximum relative 
error during one period $f_{err}(H)$:

\begin{equation}
\label{f-err}
f_{err} =  \sup_{t \in [0,2\pi/\omega]} \sqrt{
\frac{(X_{ser}-X_{num})^2 + (dX_{ser}/dt-dX_{num}/dt)^2}
{X_{ser}^2+(dX_{ser}/dt)^2}}
\end{equation}

This function indicates a maximum relative deviation in phase space
between series (\ref{res-duffing}) $X_{ser}$ and its symbolically 
evaluated derivation on the one hand and numerical solutions of
(\ref{eq-duffing}) $X_{num}, dX_{num}/dt$ on the another hand 
as a function of 
full energy $H$. We have:

$$
\begin{array}{ll}
f_{err}(H=0.1) &\simeq 1.8 \times 10^{-8}\\
f_{err}(H=0.125) &\simeq 7.4 \times 10^{-7}\\
f_{err}(H_{max}=0.163) &\simeq 7.4 \times 10^{-5}
\end{array}
$$

You can see that riched precision can be used for practical goals in
a physical diapason of energy. A use of a ready finite formulae can
be sometimes more preferable at a real time calculations than a 
numerical solution of differential equations. Of course form
(\ref{res-duffing}) must be improved for a real usage by standard
methods of prepearing of series for numerical tabulation.


\subsection{van der pol's equation}

   This equation origins from a problem of vibrations in electronic
circuits:

\begin{equation}
\label{eq-vdp}
\frac{d^{2}X}{dt^2} \ = \ - X + (\varepsilon^2 - X^2) \ \frac{dX}{dt}
\end{equation}

By linear complex changing variables (\ref{line-norm})
it can be rewritten in the diagonalized form:

$$
\begin{array}{lcrl}
\frac{\mathstrut dY_1}{\mathstrut dt}& = &  i &\!\! Y_1 + 
   {1 \over 2} \ (Y_1-Y_2) \ [\varepsilon^2 - (Y_1 +Y_2)^2]\\
\frac{\mathstrut dY_2}{\mathstrut dt}& = & -i &\!\! Y_2 + 
   {1 \over 2} \ (Y_2-Y_1) \ [\varepsilon^2 - (Y_1 +Y_2)^2]\\
\frac{\mathstrut d\varepsilon}{\mathstrut dt}& = &0
\end{array}
$$

\noindent
As for Duffings` case these equations have complex conjugated 
coefficients at terms with exchange $Y_1 \leftrightarrow Y_2$.

We consider here the  parameter $\varepsilon$ as a new variable.
Such trick allows us to free eigenvalues from a parameter 
dependence. It is necessary to be careful under this operation. We
can produce the trick only if the expansion of the expression $1/
(\overline{Q}_i,\overline{\lambda}(\varepsilon))$ is correct
at the point $\varepsilon = 0$ for all $\overline{Q}_i$ 
used in the summation in (\ref{n-trans}). (See the footnote on 
the correspondent page).

And as for Duffing's equation sums at the right sides of the 
normal form  will include only terms where $p_1 = p_2$. The third 
"additional" equation will be not changed.
So a difference between (\ref{n-duffing}) and the system presented
below is only in polynomial 
dependence of the right sides on $\varepsilon$:

\begin{equation}
\label{n-vdp}
\begin{array}{lcrl}
\frac{\mathstrut dZ_1}{\mathstrut dt}& = &  i &\!\! Z_1 
+ Z_1\ \sum_{k=0,1,\ldots}\ (g_{1,1,1,2k} (Z_1 Z_2) \varepsilon^{2k} + 
        g_{1,2,2,2k}\ (Z_1 Z_2)^2 \varepsilon^{2k}
         + \ldots)\\
\frac{\mathstrut dZ_2}{\mathstrut dt}& = & -i &\!\! Z_2
+ Z_2\ \sum_{k=0,1,\ldots}\ (g_{2,1,1,2k} (Z_1 Z_2) \varepsilon^{2k} + 
        g_{2,2,2,2k}(\varepsilon) (Z_1 Z_2)^2 \varepsilon^{2k} + \ldots)\\
\frac{\mathstrut d\varepsilon}{\mathstrut dt}& = &0&
\end{array}
\end{equation}

The convergent condition (\ref{cond-2}) will have now the form:

\begin{equation}
\label{cond-vdp}
Re[\sum_{j,k=0,1,\ldots} g_{1,j,j,2k} \ \varepsilon^{2k} \ 
(Z_1 Z_2)^j] = 0,\qquad 
g_{1,0,0,0} \stackrel{\rm def}{=} 0
\end{equation}

Contrary to Duffing's equation the above is not satisfied automatically,
but it may be rewrite in the form:

$$
a U + b V + c U^2 +d U V + e V^2 + \ldots = 0
$$

\noindent 
with $U = Z_1 Z_2, V = \varepsilon^2$ and $a \neq 0$.
This is an experimental fact and it allows us 
to get the unique (because of the implicit function theorem)
solution of (\ref{cond-vdp}) in the form:

$$
Z_1 Z_2 = \sum_{k=1,2,\ldots} q_k \varepsilon^{2k}
$$

It is easy to see that if the above is satisfied then $Z_1 Z_2$ is a
constant in time. So we can continue the evaluations as for Duffing's
case, but now constancys $c_1, c_2$ are not free (this is a case of 
a"limit circle"):  


$$ 
\begin{array}{c}
Z_1(t) = c_1 e^{+ i \omega(c_1 c_2) t}, \quad
Z_2(t) = c_2 e^{- i \omega(c_1 c_2) t} \\
c_1 c_2 = \sum_{k=1,2,\ldots} q_k \varepsilon^{2k}
\end{array}
$$


After we create the approximation to the origin equation 
(\ref{eq-vdp}) by substitution of the found above $Z_i$ into $Y_i$ with 
the help of (\ref{n-trans})
and next in $X$. If we choose complex conjugate values for 
$c_1, c_2$ we also for Duffing's equation will have an approximation 
of the real solution in the form of truncated Fourier series.

But  form (\ref{eq-vdp})
is not well for a graphic representation because  \footnote{Really, 
you can see from the above that $Z_1 Z_2 \sim \varepsilon^2$ and 
$X \sim Y_i \sim Z_i$.}
a vibration 
amplitude $X$ is proportional to the value of $\varepsilon$ and
in a picture we will have different scales for different
values of the parameter. We eliminate this problem by introducing
for the end result a new variable $\chi = X/\varepsilon$. The equation
for this variable will have the form:

\begin{equation}
\label{eq-newvdp}
\frac{d^{2}\chi}{dt^2} \ = \ - \chi + 
\varepsilon^2 \ (1 - \chi^2) \ \frac{d\chi}{dt}
\end{equation}

We did not choose this form for the investigation because in form
(\ref{eq-vdp}) we choose variables $X \sim \varepsilon$ as having
the same weight near zero. Actually we could use form (\ref{eq-newvdp})
but the result would have smaller precision at the same order of the
corresponding normal form. For (\ref{eq-newvdp}) we will have (we
display it here only till the $16^{th}$ order):


\begin{equation}
\label{res-newvdp}
\begin{array}{ll}

\omega =& 1
-{\mathstrut 1 \over \mathstrut  16} \varepsilon^4
+{\mathstrut 17 \over \mathstrut  3072} \varepsilon^8
+{\mathstrut 35 \over \mathstrut  884736} \varepsilon^{12}
-{\mathstrut 678899 \over \mathstrut  5096079360} 
\varepsilon^{16} + \ldots\\

&\\

% u=x/ep, u''-\varepsilon^2(1-u^2)u'+u=0$

\chi = &\ \ \cos(\omega t) (2
+{\mathstrut 1 \over \mathstrut  64} \varepsilon^4
-{\mathstrut 23 \over \mathstrut  49152} \varepsilon^8
-{\mathstrut 51619 \over \mathstrut  169869312} \varepsilon^{12}
+{\mathstrut 948555443 \over \mathstrut  19568944742400} 
\varepsilon^{16})\\
&+\cos(3\omega t) \varepsilon^4 (
-{\mathstrut 3 \over \mathstrut  32}
+{\mathstrut 101 \over \mathstrut  12288} \varepsilon^4
+{\mathstrut 24061 \over \mathstrut  28311552} \varepsilon^8
-{\mathstrut 279818087 \over \mathstrut  815372697600} 
\varepsilon^{12})\\
&+\cos(5\omega t) \varepsilon^4 (
-{\mathstrut 5 \over \mathstrut  96}
+{\mathstrut 1865 \over \mathstrut  110592} \varepsilon^4
-{\mathstrut 328835 \over \mathstrut  254803968} \varepsilon^8
-{\mathstrut 111998015 \over \mathstrut  293534171136} 
\varepsilon^{12})\\
&+\cos(7\omega t) \varepsilon^8 ({\mathstrut 1379 \over \mathstrut  110592}
-{\mathstrut 10923199 \over \mathstrut  3185049600} \varepsilon^4
+{\mathstrut 21049213549 \over \mathstrut  183458856960000} \varepsilon^8)\\
&+\cos(9\omega t) \varepsilon^8 ({\mathstrut 61 \over \mathstrut  20480}
-{\mathstrut 1769369 \over \mathstrut  589824000} \varepsilon^4
+{\mathstrut 161113663733 \over \mathstrut  237817036800000} \varepsilon^8)\\
&+\cos(11\omega t) \varepsilon^{12} (
-{\mathstrut 409871 \over \mathstrut  331776000}
+{\mathstrut 1359229760383 \over \mathstrut  1872809164800000} \varepsilon^4)\\
&+\cos(13\omega t) \varepsilon^{12} (
-{\mathstrut 715247 \over \mathstrut  3715891200}
+{\mathstrut 2076538440769 \over \mathstrut  5243865661440000} \varepsilon^4)\\
&+\cos(15\omega t) \varepsilon^{16} 
({\mathstrut 526426361 \over \mathstrut  4661213921280})\\
&+\cos(17\omega t) \varepsilon^{16} 
({\mathstrut 392636471 \over \mathstrut  29964946636800})\\
&+\sin(3\omega t) \varepsilon^2 (
-{\mathstrut 1 \over \mathstrut  4}
+{\mathstrut 15 \over \mathstrut  512} \varepsilon^4
-{\mathstrut 779 \over \mathstrut  1179648} \varepsilon^8
-{\mathstrut 4538017 \over \mathstrut  6794772480} 
\varepsilon^{12})\\
&+\sin(5\omega t) \varepsilon^6 ({\mathstrut 85 \over \mathstrut  2304}
-{\mathstrut 8095 \over \mathstrut  1327104} \varepsilon^4
-{\mathstrut 1252495 \over \mathstrut  6115295232} \varepsilon^8)\\
&+\sin(7\omega t) \varepsilon^6 ({\mathstrut 7 \over \mathstrut  576}
-{\mathstrut 99967 \over \mathstrut  13271040} \varepsilon^4
+{\mathstrut 415949513 \over \mathstrut  382205952000} \varepsilon^8)\\
&+\sin(9\omega t) \varepsilon^{10} (
-{\mathstrut 9791 \over \mathstrut  2457600}
+{\mathstrut 117258703 \over \mathstrut  70778880000} \varepsilon^4)\\
&+\sin(11\omega t) \varepsilon^{10} (
-{\mathstrut 5533 \over \mathstrut  7372800}
+{\mathstrut 1657839733 \over \mathstrut  1486356480000} \varepsilon^4)\\
&+\sin(13\omega t) \varepsilon^{14} 
({\mathstrut 21731177 \over \mathstrut  57802752000})\\
&+\sin(15\omega t) \varepsilon^{14} 
({\mathstrut 138697 \over \mathstrut  2774532096}) + \ldots

\end{array}
\end{equation}

   The calculation by the NORT package till $28^{\mbox th}$ order in 
$\varepsilon$ took 21 minutes on the same computer as for Duffing's case. 
We had 113 terms for  
each sums in the normal form (\ref{n-vdp}) and 1232 terms for
the each normalizing transformation from $Z_i$ to $Y_i$. The 
expression for a frequency has 9 terms and the expression for $X(t)$ 
has 114 terms. Note, that the power series for a frequency of van der 
Pol's equation has been calculated till $164^{th}$ order in $\varepsilon$ in
(Andersen, Geer, 1983).

The comparison of this result with a numerical one in the terms 
(\ref{f-err}) gives:

$$
\begin{array}{ll}
f_{err}(\varepsilon^2=0.5) &\simeq 8.9 \times 10^{-9}\\
f_{err}(\varepsilon^2=0.75) &\simeq 4.1 \times 10^{-6}\\
f_{err}(\varepsilon^2=1.0) &\simeq 2.9 \times 10^{-4}
\end{array}
$$

\noindent
Results of the tabulation of a full (in 28 order) variant of 
(\ref{res-newvdp}) for the above values of $\varepsilon$ are presented in 
Fig.\ref{fig-newvdp}.

\begin{figure}
\label{fig-newvdp}
\vspace{1cm}
\fbox{$\frac{d^{2}\chi}{dt^2} \ = \ - \chi + 
       \varepsilon^2 \ (1 - \chi^2) \ \frac{d\chi}{dt}$}

% GNUPLOT: LaTeX picture
\setlength{\unitlength}{0.240900pt}
\ifx\plotpoint\undefined\newsavebox{\plotpoint}\fi
\sbox{\plotpoint}{\rule[-0.500pt]{1.000pt}{1.000pt}}%
\begin{picture}(1500,900)(0,0)
\font\gnuplot=cmr10 at 10pt
\gnuplot
\sbox{\plotpoint}{\rule[-0.500pt]{1.000pt}{1.000pt}}%
\put(220.0,495.0){\rule[-0.500pt]{292.934pt}{1.000pt}}
\put(828.0,113.0){\rule[-0.500pt]{1.000pt}{184.048pt}}
\put(220.0,113.0){\rule[-0.500pt]{4.818pt}{1.000pt}}
\put(198,113){\makebox(0,0)[r]{-3}}
\put(1416.0,113.0){\rule[-0.500pt]{4.818pt}{1.000pt}}
\put(220.0,240.0){\rule[-0.500pt]{4.818pt}{1.000pt}}
\put(198,240){\makebox(0,0)[r]{-2}}
\put(1416.0,240.0){\rule[-0.500pt]{4.818pt}{1.000pt}}
\put(220.0,368.0){\rule[-0.500pt]{4.818pt}{1.000pt}}
\put(198,368){\makebox(0,0)[r]{-1}}
\put(1416.0,368.0){\rule[-0.500pt]{4.818pt}{1.000pt}}
\put(220.0,495.0){\rule[-0.500pt]{4.818pt}{1.000pt}}
\put(198,495){\makebox(0,0)[r]{0}}
\put(1416.0,495.0){\rule[-0.500pt]{4.818pt}{1.000pt}}
\put(220.0,622.0){\rule[-0.500pt]{4.818pt}{1.000pt}}
\put(198,622){\makebox(0,0)[r]{1}}
\put(1416.0,622.0){\rule[-0.500pt]{4.818pt}{1.000pt}}
\put(220.0,750.0){\rule[-0.500pt]{4.818pt}{1.000pt}}
\put(198,750){\makebox(0,0)[r]{2}}
\put(1416.0,750.0){\rule[-0.500pt]{4.818pt}{1.000pt}}
\put(220.0,877.0){\rule[-0.500pt]{4.818pt}{1.000pt}}
\put(198,877){\makebox(0,0)[r]{3}}
\put(1416.0,877.0){\rule[-0.500pt]{4.818pt}{1.000pt}}
\put(220.0,113.0){\rule[-0.500pt]{1.000pt}{4.818pt}}
\put(220,68){\makebox(0,0){-2.5}}
\put(220.0,857.0){\rule[-0.500pt]{1.000pt}{4.818pt}}
\put(342.0,113.0){\rule[-0.500pt]{1.000pt}{4.818pt}}
\put(342,68){\makebox(0,0){-2}}
\put(342.0,857.0){\rule[-0.500pt]{1.000pt}{4.818pt}}
\put(463.0,113.0){\rule[-0.500pt]{1.000pt}{4.818pt}}
\put(463,68){\makebox(0,0){-1.5}}
\put(463.0,857.0){\rule[-0.500pt]{1.000pt}{4.818pt}}
\put(585.0,113.0){\rule[-0.500pt]{1.000pt}{4.818pt}}
\put(585,68){\makebox(0,0){-1}}
\put(585.0,857.0){\rule[-0.500pt]{1.000pt}{4.818pt}}
\put(706.0,113.0){\rule[-0.500pt]{1.000pt}{4.818pt}}
\put(706,68){\makebox(0,0){-0.5}}
\put(706.0,857.0){\rule[-0.500pt]{1.000pt}{4.818pt}}
\put(828.0,113.0){\rule[-0.500pt]{1.000pt}{4.818pt}}
\put(828,68){\makebox(0,0){0}}
\put(828.0,857.0){\rule[-0.500pt]{1.000pt}{4.818pt}}
\put(950.0,113.0){\rule[-0.500pt]{1.000pt}{4.818pt}}
\put(950,68){\makebox(0,0){0.5}}
\put(950.0,857.0){\rule[-0.500pt]{1.000pt}{4.818pt}}
\put(1071.0,113.0){\rule[-0.500pt]{1.000pt}{4.818pt}}
\put(1071,68){\makebox(0,0){1}}
\put(1071.0,857.0){\rule[-0.500pt]{1.000pt}{4.818pt}}
\put(1193.0,113.0){\rule[-0.500pt]{1.000pt}{4.818pt}}
\put(1193,68){\makebox(0,0){1.5}}
\put(1193.0,857.0){\rule[-0.500pt]{1.000pt}{4.818pt}}
\put(1314.0,113.0){\rule[-0.500pt]{1.000pt}{4.818pt}}
\put(1314,68){\makebox(0,0){2}}
\put(1314.0,857.0){\rule[-0.500pt]{1.000pt}{4.818pt}}
\put(1436.0,113.0){\rule[-0.500pt]{1.000pt}{4.818pt}}
\put(1436,68){\makebox(0,0){2.5}}
\put(1436.0,857.0){\rule[-0.500pt]{1.000pt}{4.818pt}}
\put(220.0,113.0){\rule[-0.500pt]{292.934pt}{1.000pt}}
\put(1436.0,113.0){\rule[-0.500pt]{1.000pt}{184.048pt}}
\put(220.0,877.0){\rule[-0.500pt]{292.934pt}{1.000pt}}
\put(133,495){\makebox(0,0){$\dot{\chi}$}}
\put(828,23){\makebox(0,0){$\chi$}}
\put(220.0,113.0){\rule[-0.500pt]{1.000pt}{184.048pt}}
\sbox{\plotpoint}{\rule[-0.300pt]{0.600pt}{0.600pt}}%
\put(1306,812){\makebox(0,0)[r]{$\varepsilon^2 = 0.50$}}
\put(1328.0,812.0){\rule[-0.300pt]{15.899pt}{0.600pt}}
\put(1307,453){\usebox{\plotpoint}}
\multiput(1305.50,447.81)(-0.501,-0.969){7}{\rule{0.121pt}{1.250pt}}
\multiput(1305.75,450.41)(-6.000,-8.406){2}{\rule{0.600pt}{0.625pt}}
\multiput(1299.50,437.46)(-0.501,-0.808){9}{\rule{0.121pt}{1.093pt}}
\multiput(1299.75,439.73)(-7.000,-8.732){2}{\rule{0.600pt}{0.546pt}}
\multiput(1292.50,427.26)(-0.501,-0.626){11}{\rule{0.121pt}{0.900pt}}
\multiput(1292.75,429.13)(-8.000,-8.132){2}{\rule{0.600pt}{0.450pt}}
\multiput(1282.61,419.50)(-0.551,-0.501){13}{\rule{0.817pt}{0.121pt}}
\multiput(1284.30,419.75)(-8.305,-9.000){2}{\rule{0.408pt}{0.600pt}}
\multiput(1272.26,410.50)(-0.626,-0.501){11}{\rule{0.900pt}{0.121pt}}
\multiput(1274.13,410.75)(-8.132,-8.000){2}{\rule{0.450pt}{0.600pt}}
\multiput(1261.64,402.50)(-0.764,-0.501){11}{\rule{1.050pt}{0.121pt}}
\multiput(1263.82,402.75)(-9.821,-8.000){2}{\rule{0.525pt}{0.600pt}}
\multiput(1249.33,394.50)(-0.833,-0.501){11}{\rule{1.125pt}{0.121pt}}
\multiput(1251.67,394.75)(-10.665,-8.000){2}{\rule{0.563pt}{0.600pt}}
\multiput(1235.75,386.50)(-0.969,-0.501){9}{\rule{1.264pt}{0.121pt}}
\multiput(1238.38,386.75)(-10.376,-7.000){2}{\rule{0.632pt}{0.600pt}}
\multiput(1222.40,379.50)(-1.050,-0.501){9}{\rule{1.350pt}{0.121pt}}
\multiput(1225.20,379.75)(-11.198,-7.000){2}{\rule{0.675pt}{0.600pt}}
\multiput(1208.04,372.50)(-1.130,-0.501){9}{\rule{1.436pt}{0.121pt}}
\multiput(1211.02,372.75)(-12.020,-7.000){2}{\rule{0.718pt}{0.600pt}}
\multiput(1192.68,365.50)(-1.211,-0.501){9}{\rule{1.521pt}{0.121pt}}
\multiput(1195.84,365.75)(-12.842,-7.000){2}{\rule{0.761pt}{0.600pt}}
\multiput(1176.33,358.50)(-1.291,-0.501){9}{\rule{1.607pt}{0.121pt}}
\multiput(1179.66,358.75)(-13.664,-7.000){2}{\rule{0.804pt}{0.600pt}}
\multiput(1158.97,351.50)(-1.372,-0.501){9}{\rule{1.693pt}{0.121pt}}
\multiput(1162.49,351.75)(-14.486,-7.000){2}{\rule{0.846pt}{0.600pt}}
\multiput(1140.62,344.50)(-1.452,-0.501){9}{\rule{1.779pt}{0.121pt}}
\multiput(1144.31,344.75)(-15.308,-7.000){2}{\rule{0.889pt}{0.600pt}}
\multiput(1122.46,337.50)(-1.246,-0.501){11}{\rule{1.575pt}{0.121pt}}
\multiput(1125.73,337.75)(-15.731,-8.000){2}{\rule{0.788pt}{0.600pt}}
\multiput(1102.84,329.50)(-1.384,-0.501){11}{\rule{1.725pt}{0.121pt}}
\multiput(1106.42,329.75)(-17.420,-8.000){2}{\rule{0.863pt}{0.600pt}}
\multiput(1081.84,321.50)(-1.384,-0.501){11}{\rule{1.725pt}{0.121pt}}
\multiput(1085.42,321.75)(-17.420,-8.000){2}{\rule{0.863pt}{0.600pt}}
\multiput(1060.53,313.50)(-1.453,-0.501){11}{\rule{1.800pt}{0.121pt}}
\multiput(1064.26,313.75)(-18.264,-8.000){2}{\rule{0.900pt}{0.600pt}}
\multiput(1038.74,305.50)(-1.395,-0.501){13}{\rule{1.750pt}{0.121pt}}
\multiput(1042.37,305.75)(-20.368,-9.000){2}{\rule{0.875pt}{0.600pt}}
\multiput(1014.74,296.50)(-1.395,-0.501){13}{\rule{1.750pt}{0.121pt}}
\multiput(1018.37,296.75)(-20.368,-9.000){2}{\rule{0.875pt}{0.600pt}}
\multiput(990.18,287.50)(-1.515,-0.501){13}{\rule{1.883pt}{0.121pt}}
\multiput(994.09,287.75)(-22.091,-9.000){2}{\rule{0.942pt}{0.600pt}}
\multiput(963.91,278.50)(-1.575,-0.501){13}{\rule{1.950pt}{0.121pt}}
\multiput(967.95,278.75)(-22.953,-9.000){2}{\rule{0.975pt}{0.600pt}}
\multiput(936.63,269.50)(-1.636,-0.501){13}{\rule{2.017pt}{0.121pt}}
\multiput(940.81,269.75)(-23.814,-9.000){2}{\rule{1.008pt}{0.600pt}}
\multiput(909.15,260.50)(-1.511,-0.501){15}{\rule{1.890pt}{0.121pt}}
\multiput(913.08,260.75)(-25.077,-10.000){2}{\rule{0.945pt}{0.600pt}}
\multiput(879.08,250.50)(-1.756,-0.501){13}{\rule{2.150pt}{0.121pt}}
\multiput(883.54,250.75)(-25.538,-9.000){2}{\rule{1.075pt}{0.600pt}}
\multiput(848.80,241.50)(-1.817,-0.501){13}{\rule{2.217pt}{0.121pt}}
\multiput(853.40,241.75)(-26.399,-9.000){2}{\rule{1.108pt}{0.600pt}}
\multiput(814.99,232.50)(-2.499,-0.501){9}{\rule{2.893pt}{0.121pt}}
\multiput(821.00,232.75)(-25.996,-7.000){2}{\rule{1.446pt}{0.600pt}}
\multiput(782.64,225.50)(-2.579,-0.501){9}{\rule{2.979pt}{0.121pt}}
\multiput(788.82,225.75)(-26.818,-7.000){2}{\rule{1.489pt}{0.600pt}}
\multiput(744.44,218.50)(-4.110,-0.502){5}{\rule{4.230pt}{0.121pt}}
\multiput(753.22,218.75)(-25.220,-5.000){2}{\rule{2.115pt}{0.600pt}}
\put(694,212.25){\rule{6.950pt}{0.600pt}}
\multiput(713.57,213.75)(-19.575,-3.000){2}{\rule{3.475pt}{0.600pt}}
\put(625,211.75){\rule{8.191pt}{0.600pt}}
\multiput(642.00,210.75)(-17.000,2.000){2}{\rule{4.095pt}{0.600pt}}
\multiput(610.26,214.99)(-3.211,0.501){7}{\rule{3.550pt}{0.121pt}}
\multiput(617.63,212.75)(-26.632,6.000){2}{\rule{1.775pt}{0.600pt}}
\multiput(581.24,220.99)(-1.937,0.501){13}{\rule{2.350pt}{0.121pt}}
\multiput(586.12,218.75)(-28.122,9.000){2}{\rule{1.175pt}{0.600pt}}
\multiput(551.25,230.00)(-1.262,0.500){21}{\rule{1.627pt}{0.121pt}}
\multiput(554.62,227.75)(-28.623,13.000){2}{\rule{0.813pt}{0.600pt}}
\multiput(520.71,243.00)(-0.950,0.500){27}{\rule{1.275pt}{0.121pt}}
\multiput(523.35,240.75)(-27.354,16.000){2}{\rule{0.638pt}{0.600pt}}
\multiput(491.71,259.00)(-0.741,0.500){33}{\rule{1.034pt}{0.121pt}}
\multiput(493.85,256.75)(-25.853,19.000){2}{\rule{0.517pt}{0.600pt}}
\multiput(464.41,278.00)(-0.595,0.500){37}{\rule{0.864pt}{0.121pt}}
\multiput(466.21,275.75)(-23.206,21.000){2}{\rule{0.432pt}{0.600pt}}
\multiput(439.89,299.00)(-0.497,0.500){41}{\rule{0.750pt}{0.121pt}}
\multiput(441.44,296.75)(-21.443,23.000){2}{\rule{0.375pt}{0.600pt}}
\multiput(418.50,321.00)(-0.500,0.632){33}{\rule{0.121pt}{0.908pt}}
\multiput(418.75,321.00)(-19.000,22.116){2}{\rule{0.600pt}{0.454pt}}
\multiput(399.50,345.00)(-0.500,0.740){29}{\rule{0.121pt}{1.032pt}}
\multiput(399.75,345.00)(-17.000,22.857){2}{\rule{0.600pt}{0.516pt}}
\multiput(382.50,370.00)(-0.500,0.869){23}{\rule{0.121pt}{1.179pt}}
\multiput(382.75,370.00)(-14.000,21.554){2}{\rule{0.600pt}{0.589pt}}
\multiput(368.50,394.00)(-0.501,1.073){17}{\rule{0.121pt}{1.405pt}}
\multiput(368.75,394.00)(-11.000,20.085){2}{\rule{0.600pt}{0.702pt}}
\multiput(357.50,417.00)(-0.501,1.384){11}{\rule{0.121pt}{1.725pt}}
\multiput(357.75,417.00)(-8.000,17.420){2}{\rule{0.600pt}{0.863pt}}
\multiput(349.50,438.00)(-0.501,1.943){7}{\rule{0.121pt}{2.250pt}}
\multiput(349.75,438.00)(-6.000,16.330){2}{\rule{0.600pt}{1.125pt}}
\put(342.25,459){\rule{0.600pt}{3.750pt}}
\multiput(343.75,459.00)(-3.000,10.217){2}{\rule{0.600pt}{1.875pt}}
\put(340.25,477){\rule{0.600pt}{4.095pt}}
\multiput(340.75,477.00)(-1.000,8.500){2}{\rule{0.600pt}{2.048pt}}
\put(340.25,494){\rule{0.600pt}{3.854pt}}
\multiput(339.75,494.00)(1.000,8.000){2}{\rule{0.600pt}{1.927pt}}
\put(342.25,510){\rule{0.600pt}{2.950pt}}
\multiput(340.75,510.00)(3.000,7.877){2}{\rule{0.600pt}{1.475pt}}
\multiput(345.99,524.00)(0.503,2.141){3}{\rule{0.121pt}{2.100pt}}
\multiput(343.75,524.00)(4.000,8.641){2}{\rule{0.600pt}{1.050pt}}
\multiput(349.99,537.00)(0.501,0.969){7}{\rule{0.121pt}{1.250pt}}
\multiput(347.75,537.00)(6.000,8.406){2}{\rule{0.600pt}{0.625pt}}
\multiput(355.99,548.00)(0.501,0.808){9}{\rule{0.121pt}{1.093pt}}
\multiput(353.75,548.00)(7.000,8.732){2}{\rule{0.600pt}{0.546pt}}
\multiput(362.99,559.00)(0.501,0.626){11}{\rule{0.121pt}{0.900pt}}
\multiput(360.75,559.00)(8.000,8.132){2}{\rule{0.600pt}{0.450pt}}
\multiput(370.00,569.99)(0.551,0.501){13}{\rule{0.817pt}{0.121pt}}
\multiput(370.00,567.75)(8.305,9.000){2}{\rule{0.408pt}{0.600pt}}
\multiput(380.00,578.99)(0.626,0.501){11}{\rule{0.900pt}{0.121pt}}
\multiput(380.00,576.75)(8.132,8.000){2}{\rule{0.450pt}{0.600pt}}
\multiput(390.00,586.99)(0.764,0.501){11}{\rule{1.050pt}{0.121pt}}
\multiput(390.00,584.75)(9.821,8.000){2}{\rule{0.525pt}{0.600pt}}
\multiput(402.00,594.99)(0.833,0.501){11}{\rule{1.125pt}{0.121pt}}
\multiput(402.00,592.75)(10.665,8.000){2}{\rule{0.563pt}{0.600pt}}
\multiput(415.00,602.99)(0.969,0.501){9}{\rule{1.264pt}{0.121pt}}
\multiput(415.00,600.75)(10.376,7.000){2}{\rule{0.632pt}{0.600pt}}
\multiput(428.00,609.99)(1.050,0.501){9}{\rule{1.350pt}{0.121pt}}
\multiput(428.00,607.75)(11.198,7.000){2}{\rule{0.675pt}{0.600pt}}
\multiput(442.00,616.99)(1.130,0.501){9}{\rule{1.436pt}{0.121pt}}
\multiput(442.00,614.75)(12.020,7.000){2}{\rule{0.718pt}{0.600pt}}
\multiput(457.00,623.99)(1.211,0.501){9}{\rule{1.521pt}{0.121pt}}
\multiput(457.00,621.75)(12.842,7.000){2}{\rule{0.761pt}{0.600pt}}
\multiput(473.00,630.99)(1.291,0.501){9}{\rule{1.607pt}{0.121pt}}
\multiput(473.00,628.75)(13.664,7.000){2}{\rule{0.804pt}{0.600pt}}
\multiput(490.00,637.99)(1.372,0.501){9}{\rule{1.693pt}{0.121pt}}
\multiput(490.00,635.75)(14.486,7.000){2}{\rule{0.846pt}{0.600pt}}
\multiput(508.00,644.99)(1.452,0.501){9}{\rule{1.779pt}{0.121pt}}
\multiput(508.00,642.75)(15.308,7.000){2}{\rule{0.889pt}{0.600pt}}
\multiput(527.00,651.99)(1.246,0.501){11}{\rule{1.575pt}{0.121pt}}
\multiput(527.00,649.75)(15.731,8.000){2}{\rule{0.788pt}{0.600pt}}
\multiput(546.00,659.99)(1.384,0.501){11}{\rule{1.725pt}{0.121pt}}
\multiput(546.00,657.75)(17.420,8.000){2}{\rule{0.863pt}{0.600pt}}
\multiput(567.00,667.99)(1.384,0.501){11}{\rule{1.725pt}{0.121pt}}
\multiput(567.00,665.75)(17.420,8.000){2}{\rule{0.863pt}{0.600pt}}
\multiput(588.00,675.99)(1.453,0.501){11}{\rule{1.800pt}{0.121pt}}
\multiput(588.00,673.75)(18.264,8.000){2}{\rule{0.900pt}{0.600pt}}
\multiput(610.00,683.99)(1.395,0.501){13}{\rule{1.750pt}{0.121pt}}
\multiput(610.00,681.75)(20.368,9.000){2}{\rule{0.875pt}{0.600pt}}
\multiput(634.00,692.99)(1.395,0.501){13}{\rule{1.750pt}{0.121pt}}
\multiput(634.00,690.75)(20.368,9.000){2}{\rule{0.875pt}{0.600pt}}
\multiput(658.00,701.99)(1.515,0.501){13}{\rule{1.883pt}{0.121pt}}
\multiput(658.00,699.75)(22.091,9.000){2}{\rule{0.942pt}{0.600pt}}
\multiput(684.00,710.99)(1.575,0.501){13}{\rule{1.950pt}{0.121pt}}
\multiput(684.00,708.75)(22.953,9.000){2}{\rule{0.975pt}{0.600pt}}
\multiput(711.00,719.99)(1.636,0.501){13}{\rule{2.017pt}{0.121pt}}
\multiput(711.00,717.75)(23.814,9.000){2}{\rule{1.008pt}{0.600pt}}
\multiput(739.00,729.00)(1.511,0.501){15}{\rule{1.890pt}{0.121pt}}
\multiput(739.00,726.75)(25.077,10.000){2}{\rule{0.945pt}{0.600pt}}
\multiput(768.00,738.99)(1.756,0.501){13}{\rule{2.150pt}{0.121pt}}
\multiput(768.00,736.75)(25.538,9.000){2}{\rule{1.075pt}{0.600pt}}
\multiput(798.00,747.99)(1.817,0.501){13}{\rule{2.217pt}{0.121pt}}
\multiput(798.00,745.75)(26.399,9.000){2}{\rule{1.108pt}{0.600pt}}
\multiput(829.00,756.99)(2.499,0.501){9}{\rule{2.893pt}{0.121pt}}
\multiput(829.00,754.75)(25.996,7.000){2}{\rule{1.446pt}{0.600pt}}
\multiput(861.00,763.99)(2.579,0.501){9}{\rule{2.979pt}{0.121pt}}
\multiput(861.00,761.75)(26.818,7.000){2}{\rule{1.489pt}{0.600pt}}
\multiput(894.00,770.99)(4.110,0.502){5}{\rule{4.230pt}{0.121pt}}
\multiput(894.00,768.75)(25.220,5.000){2}{\rule{2.115pt}{0.600pt}}
\put(928,775.25){\rule{6.950pt}{0.600pt}}
\multiput(928.00,773.75)(19.575,3.000){2}{\rule{3.475pt}{0.600pt}}
\put(659.0,212.0){\rule[-0.300pt]{8.431pt}{0.600pt}}
\put(997,775.75){\rule{8.191pt}{0.600pt}}
\multiput(997.00,776.75)(17.000,-2.000){2}{\rule{4.095pt}{0.600pt}}
\multiput(1031.00,774.50)(3.211,-0.501){7}{\rule{3.550pt}{0.121pt}}
\multiput(1031.00,774.75)(26.632,-6.000){2}{\rule{1.775pt}{0.600pt}}
\multiput(1065.00,768.50)(1.937,-0.501){13}{\rule{2.350pt}{0.121pt}}
\multiput(1065.00,768.75)(28.122,-9.000){2}{\rule{1.175pt}{0.600pt}}
\multiput(1098.00,759.50)(1.262,-0.500){21}{\rule{1.627pt}{0.121pt}}
\multiput(1098.00,759.75)(28.623,-13.000){2}{\rule{0.813pt}{0.600pt}}
\multiput(1130.00,746.50)(0.950,-0.500){27}{\rule{1.275pt}{0.121pt}}
\multiput(1130.00,746.75)(27.354,-16.000){2}{\rule{0.638pt}{0.600pt}}
\multiput(1160.00,730.50)(0.741,-0.500){33}{\rule{1.034pt}{0.121pt}}
\multiput(1160.00,730.75)(25.853,-19.000){2}{\rule{0.517pt}{0.600pt}}
\multiput(1188.00,711.50)(0.595,-0.500){37}{\rule{0.864pt}{0.121pt}}
\multiput(1188.00,711.75)(23.206,-21.000){2}{\rule{0.432pt}{0.600pt}}
\multiput(1213.00,690.50)(0.497,-0.500){41}{\rule{0.750pt}{0.121pt}}
\multiput(1213.00,690.75)(21.443,-23.000){2}{\rule{0.375pt}{0.600pt}}
\multiput(1237.00,665.23)(0.500,-0.632){33}{\rule{0.121pt}{0.908pt}}
\multiput(1234.75,667.12)(19.000,-22.116){2}{\rule{0.600pt}{0.454pt}}
\multiput(1256.00,640.71)(0.500,-0.740){29}{\rule{0.121pt}{1.032pt}}
\multiput(1253.75,642.86)(17.000,-22.857){2}{\rule{0.600pt}{0.516pt}}
\multiput(1273.00,615.11)(0.500,-0.869){23}{\rule{0.121pt}{1.179pt}}
\multiput(1270.75,617.55)(14.000,-21.554){2}{\rule{0.600pt}{0.589pt}}
\multiput(1287.00,590.17)(0.501,-1.073){17}{\rule{0.121pt}{1.405pt}}
\multiput(1284.75,593.08)(11.000,-20.085){2}{\rule{0.600pt}{0.702pt}}
\multiput(1297.99,565.84)(0.501,-1.384){11}{\rule{0.121pt}{1.725pt}}
\multiput(1295.75,569.42)(8.000,-17.420){2}{\rule{0.600pt}{0.863pt}}
\multiput(1305.99,542.66)(0.501,-1.943){7}{\rule{0.121pt}{2.250pt}}
\multiput(1303.75,547.33)(6.000,-16.330){2}{\rule{0.600pt}{1.125pt}}
\put(1311.25,513){\rule{0.600pt}{3.750pt}}
\multiput(1309.75,523.22)(3.000,-10.217){2}{\rule{0.600pt}{1.875pt}}
\put(1313.25,496){\rule{0.600pt}{4.095pt}}
\multiput(1312.75,504.50)(1.000,-8.500){2}{\rule{0.600pt}{2.048pt}}
\put(1313.25,480){\rule{0.600pt}{3.854pt}}
\multiput(1313.75,488.00)(-1.000,-8.000){2}{\rule{0.600pt}{1.927pt}}
\put(1311.25,466){\rule{0.600pt}{2.950pt}}
\multiput(1312.75,473.88)(-3.000,-7.877){2}{\rule{0.600pt}{1.475pt}}
\multiput(1309.50,457.28)(-0.503,-2.141){3}{\rule{0.121pt}{2.100pt}}
\multiput(1309.75,461.64)(-4.000,-8.641){2}{\rule{0.600pt}{1.050pt}}
\put(962.0,778.0){\rule[-0.300pt]{8.431pt}{0.600pt}}
\sbox{\plotpoint}{\rule[-0.175pt]{0.350pt}{0.350pt}}%
\put(1306,767){\makebox(0,0)[r]{$\varepsilon^2 = 0.75$}}
\put(1328.0,767.0){\rule[-0.175pt]{15.899pt}{0.350pt}}
\put(1299,441){\usebox{\plotpoint}}
\multiput(1298.02,438.98)(-0.504,-0.598){11}{\rule{0.121pt}{0.488pt}}
\multiput(1298.27,439.99)(-7.000,-6.988){2}{\rule{0.350pt}{0.244pt}}
\multiput(1290.18,432.02)(-0.515,-0.504){13}{\rule{0.438pt}{0.121pt}}
\multiput(1291.09,432.27)(-7.092,-8.000){2}{\rule{0.219pt}{0.350pt}}
\multiput(1281.46,424.02)(-0.807,-0.505){9}{\rule{0.612pt}{0.122pt}}
\multiput(1282.73,424.27)(-7.729,-6.000){2}{\rule{0.306pt}{0.350pt}}
\multiput(1272.22,418.02)(-0.902,-0.505){9}{\rule{0.671pt}{0.122pt}}
\multiput(1273.61,418.27)(-8.608,-6.000){2}{\rule{0.335pt}{0.350pt}}
\multiput(1261.97,412.02)(-0.997,-0.505){9}{\rule{0.729pt}{0.122pt}}
\multiput(1263.49,412.27)(-9.487,-6.000){2}{\rule{0.365pt}{0.350pt}}
\multiput(1250.97,406.02)(-0.997,-0.505){9}{\rule{0.729pt}{0.122pt}}
\multiput(1252.49,406.27)(-9.487,-6.000){2}{\rule{0.365pt}{0.350pt}}
\multiput(1239.15,400.02)(-1.358,-0.507){7}{\rule{0.928pt}{0.122pt}}
\multiput(1241.07,400.27)(-10.075,-5.000){2}{\rule{0.464pt}{0.350pt}}
\multiput(1227.15,395.02)(-1.358,-0.507){7}{\rule{0.928pt}{0.122pt}}
\multiput(1229.07,395.27)(-10.075,-5.000){2}{\rule{0.464pt}{0.350pt}}
\multiput(1215.25,390.02)(-1.281,-0.505){9}{\rule{0.904pt}{0.122pt}}
\multiput(1217.12,390.27)(-12.123,-6.000){2}{\rule{0.452pt}{0.350pt}}
\multiput(1200.57,384.02)(-1.595,-0.507){7}{\rule{1.067pt}{0.122pt}}
\multiput(1202.78,384.27)(-11.784,-5.000){2}{\rule{0.534pt}{0.350pt}}
\multiput(1186.57,379.02)(-1.595,-0.507){7}{\rule{1.067pt}{0.122pt}}
\multiput(1188.78,379.27)(-11.784,-5.000){2}{\rule{0.534pt}{0.350pt}}
\multiput(1172.76,374.02)(-1.470,-0.505){9}{\rule{1.021pt}{0.122pt}}
\multiput(1174.88,374.27)(-13.881,-6.000){2}{\rule{0.510pt}{0.350pt}}
\multiput(1156.76,368.02)(-1.470,-0.505){9}{\rule{1.021pt}{0.122pt}}
\multiput(1158.88,368.27)(-13.881,-6.000){2}{\rule{0.510pt}{0.350pt}}
\multiput(1141.32,362.02)(-1.232,-0.504){11}{\rule{0.887pt}{0.121pt}}
\multiput(1143.16,362.27)(-14.158,-7.000){2}{\rule{0.444pt}{0.350pt}}
\multiput(1124.90,355.02)(-1.390,-0.504){11}{\rule{0.987pt}{0.121pt}}
\multiput(1126.95,355.27)(-15.950,-7.000){2}{\rule{0.494pt}{0.350pt}}
\multiput(1106.69,348.02)(-1.469,-0.504){11}{\rule{1.038pt}{0.121pt}}
\multiput(1108.85,348.27)(-16.847,-7.000){2}{\rule{0.519pt}{0.350pt}}
\multiput(1088.19,341.02)(-1.266,-0.504){13}{\rule{0.919pt}{0.121pt}}
\multiput(1090.09,341.27)(-17.093,-8.000){2}{\rule{0.459pt}{0.350pt}}
\multiput(1068.82,333.02)(-1.402,-0.504){13}{\rule{1.006pt}{0.121pt}}
\multiput(1070.91,333.27)(-18.911,-8.000){2}{\rule{0.503pt}{0.350pt}}
\multiput(1048.59,325.02)(-1.099,-0.503){17}{\rule{0.822pt}{0.121pt}}
\multiput(1050.29,325.27)(-19.293,-10.000){2}{\rule{0.411pt}{0.350pt}}
\multiput(1027.30,315.02)(-1.206,-0.503){17}{\rule{0.893pt}{0.121pt}}
\multiput(1029.15,315.27)(-21.148,-10.000){2}{\rule{0.446pt}{0.350pt}}
\multiput(1004.15,305.02)(-1.259,-0.503){17}{\rule{0.928pt}{0.121pt}}
\multiput(1006.07,305.27)(-22.075,-10.000){2}{\rule{0.464pt}{0.350pt}}
\multiput(980.49,295.02)(-1.124,-0.502){21}{\rule{0.846pt}{0.121pt}}
\multiput(982.24,295.27)(-24.244,-12.000){2}{\rule{0.423pt}{0.350pt}}
\multiput(954.49,283.02)(-1.124,-0.502){21}{\rule{0.846pt}{0.121pt}}
\multiput(956.24,283.27)(-24.244,-12.000){2}{\rule{0.423pt}{0.350pt}}
\multiput(928.13,271.02)(-1.256,-0.502){21}{\rule{0.933pt}{0.121pt}}
\multiput(930.06,271.27)(-27.063,-12.000){2}{\rule{0.467pt}{0.350pt}}
\multiput(899.52,259.02)(-1.105,-0.502){25}{\rule{0.837pt}{0.121pt}}
\multiput(901.26,259.27)(-28.262,-14.000){2}{\rule{0.419pt}{0.350pt}}
\multiput(869.06,245.02)(-1.275,-0.502){23}{\rule{0.949pt}{0.121pt}}
\multiput(871.03,245.27)(-30.030,-13.000){2}{\rule{0.475pt}{0.350pt}}
\multiput(836.95,232.02)(-1.316,-0.502){23}{\rule{0.976pt}{0.121pt}}
\multiput(838.97,232.27)(-30.974,-13.000){2}{\rule{0.488pt}{0.350pt}}
\multiput(803.40,219.02)(-1.520,-0.502){21}{\rule{1.108pt}{0.121pt}}
\multiput(805.70,219.27)(-32.700,-12.000){2}{\rule{0.554pt}{0.350pt}}
\multiput(767.41,207.02)(-1.900,-0.503){17}{\rule{1.347pt}{0.121pt}}
\multiput(770.20,207.27)(-33.203,-10.000){2}{\rule{0.674pt}{0.350pt}}
\multiput(729.74,197.02)(-2.561,-0.504){13}{\rule{1.750pt}{0.121pt}}
\multiput(733.37,197.27)(-34.368,-8.000){2}{\rule{0.875pt}{0.350pt}}
\multiput(687.59,189.02)(-4.432,-0.507){7}{\rule{2.747pt}{0.122pt}}
\multiput(693.30,189.27)(-32.297,-5.000){2}{\rule{1.374pt}{0.350pt}}
\multiput(613.44,185.47)(-3.554,0.505){9}{\rule{2.304pt}{0.122pt}}
\multiput(618.22,184.27)(-33.218,6.000){2}{\rule{1.152pt}{0.350pt}}
\multiput(579.62,191.48)(-1.812,0.502){19}{\rule{1.297pt}{0.121pt}}
\multiput(582.31,190.27)(-35.309,11.000){2}{\rule{0.648pt}{0.350pt}}
\multiput(543.65,202.48)(-1.055,0.502){31}{\rule{0.808pt}{0.121pt}}
\multiput(545.32,201.27)(-33.323,17.000){2}{\rule{0.404pt}{0.350pt}}
\multiput(509.55,219.48)(-0.727,0.501){43}{\rule{0.590pt}{0.121pt}}
\multiput(510.78,218.27)(-31.776,23.000){2}{\rule{0.295pt}{0.350pt}}
\multiput(477.02,242.48)(-0.560,0.501){51}{\rule{0.476pt}{0.121pt}}
\multiput(478.01,241.27)(-29.011,27.000){2}{\rule{0.238pt}{0.350pt}}
\multiput(448.02,269.00)(-0.501,0.583){49}{\rule{0.121pt}{0.491pt}}
\multiput(448.27,269.00)(-26.000,28.980){2}{\rule{0.350pt}{0.246pt}}
\multiput(422.02,299.00)(-0.501,0.715){41}{\rule{0.121pt}{0.581pt}}
\multiput(422.27,299.00)(-22.000,29.795){2}{\rule{0.350pt}{0.290pt}}
\multiput(400.02,330.00)(-0.501,0.831){35}{\rule{0.121pt}{0.659pt}}
\multiput(400.27,330.00)(-19.000,29.633){2}{\rule{0.350pt}{0.329pt}}
\multiput(381.02,361.00)(-0.502,1.028){27}{\rule{0.121pt}{0.787pt}}
\multiput(381.27,361.00)(-15.000,28.366){2}{\rule{0.350pt}{0.394pt}}
\multiput(366.02,391.00)(-0.502,1.282){19}{\rule{0.121pt}{0.947pt}}
\multiput(366.27,391.00)(-11.000,25.035){2}{\rule{0.350pt}{0.473pt}}
\multiput(355.02,418.00)(-0.504,1.675){13}{\rule{0.121pt}{1.181pt}}
\multiput(355.27,418.00)(-8.000,22.548){2}{\rule{0.350pt}{0.591pt}}
\multiput(347.02,443.00)(-0.507,2.541){7}{\rule{0.122pt}{1.628pt}}
\multiput(347.27,443.00)(-5.000,18.622){2}{\rule{0.350pt}{0.814pt}}
\put(341.27,465){\rule{0.350pt}{3.588pt}}
\multiput(342.27,465.00)(-2.000,12.554){2}{\rule{0.350pt}{1.794pt}}
\put(339.77,485){\rule{0.350pt}{4.095pt}}
\multiput(340.27,485.00)(-1.000,8.500){2}{\rule{0.350pt}{2.048pt}}
\put(340.27,502){\rule{0.350pt}{2.538pt}}
\multiput(339.27,502.00)(2.000,8.733){2}{\rule{0.350pt}{1.269pt}}
\multiput(342.47,516.00)(0.509,1.979){5}{\rule{0.123pt}{1.225pt}}
\multiput(341.27,516.00)(4.000,10.457){2}{\rule{0.350pt}{0.613pt}}
\multiput(346.47,529.00)(0.507,1.240){7}{\rule{0.122pt}{0.858pt}}
\multiput(345.27,529.00)(5.000,9.220){2}{\rule{0.350pt}{0.429pt}}
\multiput(351.47,540.00)(0.505,0.807){9}{\rule{0.122pt}{0.613pt}}
\multiput(350.27,540.00)(6.000,7.729){2}{\rule{0.350pt}{0.306pt}}
\multiput(357.47,549.00)(0.504,0.598){11}{\rule{0.121pt}{0.488pt}}
\multiput(356.27,549.00)(7.000,6.988){2}{\rule{0.350pt}{0.244pt}}
\multiput(364.00,557.47)(0.515,0.504){13}{\rule{0.438pt}{0.121pt}}
\multiput(364.00,556.27)(7.092,8.000){2}{\rule{0.219pt}{0.350pt}}
\multiput(372.00,565.47)(0.807,0.505){9}{\rule{0.612pt}{0.122pt}}
\multiput(372.00,564.27)(7.729,6.000){2}{\rule{0.306pt}{0.350pt}}
\multiput(381.00,571.47)(0.902,0.505){9}{\rule{0.671pt}{0.122pt}}
\multiput(381.00,570.27)(8.608,6.000){2}{\rule{0.335pt}{0.350pt}}
\multiput(391.00,577.47)(0.997,0.505){9}{\rule{0.729pt}{0.122pt}}
\multiput(391.00,576.27)(9.487,6.000){2}{\rule{0.365pt}{0.350pt}}
\multiput(402.00,583.47)(0.997,0.505){9}{\rule{0.729pt}{0.122pt}}
\multiput(402.00,582.27)(9.487,6.000){2}{\rule{0.365pt}{0.350pt}}
\multiput(413.00,589.47)(1.358,0.507){7}{\rule{0.928pt}{0.122pt}}
\multiput(413.00,588.27)(10.075,5.000){2}{\rule{0.464pt}{0.350pt}}
\multiput(425.00,594.47)(1.358,0.507){7}{\rule{0.928pt}{0.122pt}}
\multiput(425.00,593.27)(10.075,5.000){2}{\rule{0.464pt}{0.350pt}}
\multiput(437.00,599.47)(1.281,0.505){9}{\rule{0.904pt}{0.122pt}}
\multiput(437.00,598.27)(12.123,6.000){2}{\rule{0.452pt}{0.350pt}}
\multiput(451.00,605.47)(1.595,0.507){7}{\rule{1.067pt}{0.122pt}}
\multiput(451.00,604.27)(11.784,5.000){2}{\rule{0.534pt}{0.350pt}}
\multiput(465.00,610.47)(1.595,0.507){7}{\rule{1.067pt}{0.122pt}}
\multiput(465.00,609.27)(11.784,5.000){2}{\rule{0.534pt}{0.350pt}}
\multiput(479.00,615.47)(1.470,0.505){9}{\rule{1.021pt}{0.122pt}}
\multiput(479.00,614.27)(13.881,6.000){2}{\rule{0.510pt}{0.350pt}}
\multiput(495.00,621.47)(1.470,0.505){9}{\rule{1.021pt}{0.122pt}}
\multiput(495.00,620.27)(13.881,6.000){2}{\rule{0.510pt}{0.350pt}}
\multiput(511.00,627.47)(1.232,0.504){11}{\rule{0.887pt}{0.121pt}}
\multiput(511.00,626.27)(14.158,7.000){2}{\rule{0.444pt}{0.350pt}}
\multiput(527.00,634.47)(1.390,0.504){11}{\rule{0.987pt}{0.121pt}}
\multiput(527.00,633.27)(15.950,7.000){2}{\rule{0.494pt}{0.350pt}}
\multiput(545.00,641.47)(1.469,0.504){11}{\rule{1.038pt}{0.121pt}}
\multiput(545.00,640.27)(16.847,7.000){2}{\rule{0.519pt}{0.350pt}}
\multiput(564.00,648.47)(1.266,0.504){13}{\rule{0.919pt}{0.121pt}}
\multiput(564.00,647.27)(17.093,8.000){2}{\rule{0.459pt}{0.350pt}}
\multiput(583.00,656.47)(1.402,0.504){13}{\rule{1.006pt}{0.121pt}}
\multiput(583.00,655.27)(18.911,8.000){2}{\rule{0.503pt}{0.350pt}}
\multiput(604.00,664.48)(1.099,0.503){17}{\rule{0.822pt}{0.121pt}}
\multiput(604.00,663.27)(19.293,10.000){2}{\rule{0.411pt}{0.350pt}}
\multiput(625.00,674.48)(1.206,0.503){17}{\rule{0.893pt}{0.121pt}}
\multiput(625.00,673.27)(21.148,10.000){2}{\rule{0.446pt}{0.350pt}}
\multiput(648.00,684.48)(1.259,0.503){17}{\rule{0.928pt}{0.121pt}}
\multiput(648.00,683.27)(22.075,10.000){2}{\rule{0.464pt}{0.350pt}}
\multiput(672.00,694.48)(1.124,0.502){21}{\rule{0.846pt}{0.121pt}}
\multiput(672.00,693.27)(24.244,12.000){2}{\rule{0.423pt}{0.350pt}}
\multiput(698.00,706.48)(1.124,0.502){21}{\rule{0.846pt}{0.121pt}}
\multiput(698.00,705.27)(24.244,12.000){2}{\rule{0.423pt}{0.350pt}}
\multiput(724.00,718.48)(1.256,0.502){21}{\rule{0.933pt}{0.121pt}}
\multiput(724.00,717.27)(27.063,12.000){2}{\rule{0.467pt}{0.350pt}}
\multiput(753.00,730.48)(1.105,0.502){25}{\rule{0.837pt}{0.121pt}}
\multiput(753.00,729.27)(28.262,14.000){2}{\rule{0.419pt}{0.350pt}}
\multiput(783.00,744.48)(1.275,0.502){23}{\rule{0.949pt}{0.121pt}}
\multiput(783.00,743.27)(30.030,13.000){2}{\rule{0.475pt}{0.350pt}}
\multiput(815.00,757.48)(1.316,0.502){23}{\rule{0.976pt}{0.121pt}}
\multiput(815.00,756.27)(30.974,13.000){2}{\rule{0.488pt}{0.350pt}}
\multiput(848.00,770.48)(1.520,0.502){21}{\rule{1.108pt}{0.121pt}}
\multiput(848.00,769.27)(32.700,12.000){2}{\rule{0.554pt}{0.350pt}}
\multiput(883.00,782.48)(1.900,0.503){17}{\rule{1.347pt}{0.121pt}}
\multiput(883.00,781.27)(33.203,10.000){2}{\rule{0.674pt}{0.350pt}}
\multiput(919.00,792.47)(2.561,0.504){13}{\rule{1.750pt}{0.121pt}}
\multiput(919.00,791.27)(34.368,8.000){2}{\rule{0.875pt}{0.350pt}}
\multiput(957.00,800.47)(4.432,0.507){7}{\rule{2.747pt}{0.122pt}}
\multiput(957.00,799.27)(32.297,5.000){2}{\rule{1.374pt}{0.350pt}}
\put(623.0,185.0){\rule[-0.175pt]{9.154pt}{0.350pt}}
\multiput(1033.00,804.02)(3.554,-0.505){9}{\rule{2.304pt}{0.122pt}}
\multiput(1033.00,804.27)(33.218,-6.000){2}{\rule{1.152pt}{0.350pt}}
\multiput(1071.00,798.02)(1.812,-0.502){19}{\rule{1.297pt}{0.121pt}}
\multiput(1071.00,798.27)(35.309,-11.000){2}{\rule{0.648pt}{0.350pt}}
\multiput(1109.00,787.02)(1.055,-0.502){31}{\rule{0.808pt}{0.121pt}}
\multiput(1109.00,787.27)(33.323,-17.000){2}{\rule{0.404pt}{0.350pt}}
\multiput(1144.00,770.02)(0.727,-0.501){43}{\rule{0.590pt}{0.121pt}}
\multiput(1144.00,770.27)(31.776,-23.000){2}{\rule{0.295pt}{0.350pt}}
\multiput(1177.00,747.02)(0.560,-0.501){51}{\rule{0.476pt}{0.121pt}}
\multiput(1177.00,747.27)(29.011,-27.000){2}{\rule{0.238pt}{0.350pt}}
\multiput(1207.48,718.96)(0.501,-0.583){49}{\rule{0.121pt}{0.491pt}}
\multiput(1206.27,719.98)(26.000,-28.980){2}{\rule{0.350pt}{0.246pt}}
\multiput(1233.48,688.59)(0.501,-0.715){41}{\rule{0.121pt}{0.581pt}}
\multiput(1232.27,689.79)(22.000,-29.795){2}{\rule{0.350pt}{0.290pt}}
\multiput(1255.48,657.27)(0.501,-0.831){35}{\rule{0.121pt}{0.659pt}}
\multiput(1254.27,658.63)(19.000,-29.633){2}{\rule{0.350pt}{0.329pt}}
\multiput(1274.48,625.73)(0.502,-1.028){27}{\rule{0.121pt}{0.787pt}}
\multiput(1273.27,627.37)(15.000,-28.366){2}{\rule{0.350pt}{0.394pt}}
\multiput(1289.48,595.07)(0.502,-1.282){19}{\rule{0.121pt}{0.947pt}}
\multiput(1288.27,597.04)(11.000,-25.035){2}{\rule{0.350pt}{0.473pt}}
\multiput(1300.47,567.10)(0.504,-1.675){13}{\rule{0.121pt}{1.181pt}}
\multiput(1299.27,569.55)(8.000,-22.548){2}{\rule{0.350pt}{0.591pt}}
\multiput(1308.47,540.24)(0.507,-2.541){7}{\rule{0.122pt}{1.628pt}}
\multiput(1307.27,543.62)(5.000,-18.622){2}{\rule{0.350pt}{0.814pt}}
\put(1313.27,505){\rule{0.350pt}{3.588pt}}
\multiput(1312.27,517.55)(2.000,-12.554){2}{\rule{0.350pt}{1.794pt}}
\put(1314.77,488){\rule{0.350pt}{4.095pt}}
\multiput(1314.27,496.50)(1.000,-8.500){2}{\rule{0.350pt}{2.048pt}}
\put(1314.27,474){\rule{0.350pt}{2.538pt}}
\multiput(1315.27,482.73)(-2.000,-8.733){2}{\rule{0.350pt}{1.269pt}}
\multiput(1313.02,468.91)(-0.509,-1.979){5}{\rule{0.123pt}{1.225pt}}
\multiput(1313.27,471.46)(-4.000,-10.457){2}{\rule{0.350pt}{0.613pt}}
\multiput(1309.02,457.44)(-0.507,-1.240){7}{\rule{0.122pt}{0.858pt}}
\multiput(1309.27,459.22)(-5.000,-9.220){2}{\rule{0.350pt}{0.429pt}}
\multiput(1304.02,447.46)(-0.505,-0.807){9}{\rule{0.122pt}{0.613pt}}
\multiput(1304.27,448.73)(-6.000,-7.729){2}{\rule{0.350pt}{0.306pt}}
\put(995.0,805.0){\rule[-0.175pt]{9.154pt}{0.350pt}}
\sbox{\plotpoint}{\rule[-0.500pt]{1.000pt}{1.000pt}}%
\put(1306,722){\makebox(0,0)[r]{$\varepsilon^2 = 1.00$}}
\put(1328.0,722.0){\rule[-0.500pt]{15.899pt}{1.000pt}}
\put(1291,436){\usebox{\plotpoint}}
\multiput(1284.43,433.69)(-0.476,-0.462){4}{\rule{1.583pt}{0.111pt}}
\multiput(1287.71,433.92)(-4.714,-6.000){2}{\rule{0.792pt}{1.000pt}}
\put(1274,425.92){\rule{2.168pt}{1.000pt}}
\multiput(1278.50,427.92)(-4.500,-4.000){2}{\rule{1.084pt}{1.000pt}}
\multiput(1265.49,423.71)(-0.490,-0.424){2}{\rule{2.050pt}{0.102pt}}
\multiput(1269.75,423.92)(-4.745,-5.000){2}{\rule{1.025pt}{1.000pt}}
\put(1256,416.92){\rule{2.168pt}{1.000pt}}
\multiput(1260.50,418.92)(-4.500,-4.000){2}{\rule{1.084pt}{1.000pt}}
\put(1245,412.92){\rule{2.650pt}{1.000pt}}
\multiput(1250.50,414.92)(-5.500,-4.000){2}{\rule{1.325pt}{1.000pt}}
\put(1235,408.92){\rule{2.409pt}{1.000pt}}
\multiput(1240.00,410.92)(-5.000,-4.000){2}{\rule{1.204pt}{1.000pt}}
\put(1223,404.92){\rule{2.891pt}{1.000pt}}
\multiput(1229.00,406.92)(-6.000,-4.000){2}{\rule{1.445pt}{1.000pt}}
\put(1212,400.92){\rule{2.650pt}{1.000pt}}
\multiput(1217.50,402.92)(-5.500,-4.000){2}{\rule{1.325pt}{1.000pt}}
\put(1200,396.92){\rule{2.891pt}{1.000pt}}
\multiput(1206.00,398.92)(-6.000,-4.000){2}{\rule{1.445pt}{1.000pt}}
\put(1187,392.92){\rule{3.132pt}{1.000pt}}
\multiput(1193.50,394.92)(-6.500,-4.000){2}{\rule{1.566pt}{1.000pt}}
\multiput(1175.17,390.71)(-1.169,-0.424){2}{\rule{2.850pt}{0.102pt}}
\multiput(1181.08,390.92)(-7.085,-5.000){2}{\rule{1.425pt}{1.000pt}}
\multiput(1160.51,385.71)(-1.509,-0.424){2}{\rule{3.250pt}{0.102pt}}
\multiput(1167.25,385.92)(-8.254,-5.000){2}{\rule{1.625pt}{1.000pt}}
\multiput(1146.34,380.71)(-1.339,-0.424){2}{\rule{3.050pt}{0.102pt}}
\multiput(1152.67,380.92)(-7.670,-5.000){2}{\rule{1.525pt}{1.000pt}}
\multiput(1133.58,375.69)(-1.195,-0.462){4}{\rule{2.750pt}{0.111pt}}
\multiput(1139.29,375.92)(-9.292,-6.000){2}{\rule{1.375pt}{1.000pt}}
\multiput(1119.47,369.69)(-1.095,-0.475){6}{\rule{2.536pt}{0.114pt}}
\multiput(1124.74,369.92)(-10.737,-7.000){2}{\rule{1.268pt}{1.000pt}}
\multiput(1102.88,362.69)(-1.176,-0.475){6}{\rule{2.679pt}{0.114pt}}
\multiput(1108.44,362.92)(-11.440,-7.000){2}{\rule{1.339pt}{1.000pt}}
\multiput(1085.29,355.69)(-1.258,-0.475){6}{\rule{2.821pt}{0.114pt}}
\multiput(1091.14,355.92)(-12.144,-7.000){2}{\rule{1.411pt}{1.000pt}}
\multiput(1069.20,348.68)(-1.022,-0.485){10}{\rule{2.361pt}{0.117pt}}
\multiput(1074.10,348.92)(-14.099,-9.000){2}{\rule{1.181pt}{1.000pt}}
\multiput(1049.74,339.68)(-1.082,-0.485){10}{\rule{2.472pt}{0.117pt}}
\multiput(1054.87,339.92)(-14.869,-9.000){2}{\rule{1.236pt}{1.000pt}}
\multiput(1030.66,330.68)(-0.974,-0.489){14}{\rule{2.250pt}{0.118pt}}
\multiput(1035.33,330.92)(-17.330,-11.000){2}{\rule{1.125pt}{1.000pt}}
\multiput(1009.35,319.68)(-0.890,-0.491){16}{\rule{2.083pt}{0.118pt}}
\multiput(1013.68,319.92)(-17.676,-12.000){2}{\rule{1.042pt}{1.000pt}}
\multiput(986.98,307.68)(-0.940,-0.492){18}{\rule{2.173pt}{0.118pt}}
\multiput(991.49,307.92)(-20.490,-13.000){2}{\rule{1.087pt}{1.000pt}}
\multiput(962.25,294.68)(-0.908,-0.492){20}{\rule{2.107pt}{0.119pt}}
\multiput(966.63,294.92)(-21.627,-14.000){2}{\rule{1.054pt}{1.000pt}}
\multiput(936.21,280.68)(-0.915,-0.493){22}{\rule{2.117pt}{0.119pt}}
\multiput(940.61,280.92)(-23.607,-15.000){2}{\rule{1.058pt}{1.000pt}}
\multiput(908.64,265.68)(-0.865,-0.494){26}{\rule{2.015pt}{0.119pt}}
\multiput(912.82,265.92)(-25.818,-17.000){2}{\rule{1.007pt}{1.000pt}}
\multiput(878.58,248.68)(-0.873,-0.495){28}{\rule{2.028pt}{0.119pt}}
\multiput(882.79,248.92)(-27.791,-18.000){2}{\rule{1.014pt}{1.000pt}}
\multiput(846.12,230.68)(-0.930,-0.495){28}{\rule{2.139pt}{0.119pt}}
\multiput(850.56,230.92)(-29.561,-18.000){2}{\rule{1.069pt}{1.000pt}}
\multiput(811.88,212.68)(-0.961,-0.495){30}{\rule{2.197pt}{0.119pt}}
\multiput(816.44,212.92)(-32.439,-19.000){2}{\rule{1.099pt}{1.000pt}}
\multiput(772.84,193.68)(-1.211,-0.494){24}{\rule{2.688pt}{0.119pt}}
\multiput(778.42,193.92)(-33.422,-16.000){2}{\rule{1.344pt}{1.000pt}}
\multiput(731.81,177.68)(-1.463,-0.492){20}{\rule{3.179pt}{0.119pt}}
\multiput(738.40,177.92)(-34.403,-14.000){2}{\rule{1.589pt}{1.000pt}}
\multiput(683.13,163.68)(-2.459,-0.485){10}{\rule{5.028pt}{0.117pt}}
\multiput(693.56,163.92)(-32.565,-9.000){2}{\rule{2.514pt}{1.000pt}}
\put(618,153.42){\rule{10.359pt}{1.000pt}}
\multiput(639.50,154.92)(-21.500,-3.000){2}{\rule{5.179pt}{1.000pt}}
\multiput(587.21,155.84)(-4.071,0.462){4}{\rule{7.417pt}{0.111pt}}
\multiput(602.61,151.92)(-27.606,6.000){2}{\rule{3.708pt}{1.000pt}}
\multiput(563.07,161.83)(-1.307,0.494){24}{\rule{2.875pt}{0.119pt}}
\multiput(569.03,157.92)(-36.033,16.000){2}{\rule{1.438pt}{1.000pt}}
\multiput(525.22,177.83)(-0.800,0.496){40}{\rule{1.875pt}{0.120pt}}
\multiput(529.11,173.92)(-35.108,24.000){2}{\rule{0.937pt}{1.000pt}}
\multiput(488.43,201.83)(-0.534,0.497){58}{\rule{1.341pt}{0.120pt}}
\multiput(491.22,197.92)(-33.217,33.000){2}{\rule{0.670pt}{1.000pt}}
\multiput(455.68,233.00)(-0.497,0.601){54}{\rule{0.120pt}{1.476pt}}
\multiput(455.92,233.00)(-31.000,34.937){2}{\rule{1.000pt}{0.738pt}}
\multiput(424.68,271.00)(-0.497,0.758){44}{\rule{0.120pt}{1.788pt}}
\multiput(424.92,271.00)(-26.000,36.288){2}{\rule{1.000pt}{0.894pt}}
\multiput(398.68,311.00)(-0.495,0.989){32}{\rule{0.119pt}{2.250pt}}
\multiput(398.92,311.00)(-20.000,35.330){2}{\rule{1.000pt}{1.125pt}}
\multiput(378.68,351.00)(-0.494,1.146){24}{\rule{0.119pt}{2.562pt}}
\multiput(378.92,351.00)(-16.000,31.681){2}{\rule{1.000pt}{1.281pt}}
\multiput(362.68,388.00)(-0.491,1.326){16}{\rule{0.118pt}{2.917pt}}
\multiput(362.92,388.00)(-12.000,25.946){2}{\rule{1.000pt}{1.458pt}}
\multiput(350.68,420.00)(-0.481,1.845){8}{\rule{0.116pt}{3.875pt}}
\multiput(350.92,420.00)(-8.000,20.957){2}{\rule{1.000pt}{1.938pt}}
\put(340.92,449){\rule{1.000pt}{5.541pt}}
\multiput(342.92,449.00)(-4.000,11.500){2}{\rule{1.000pt}{2.770pt}}
\put(338.42,472){\rule{1.000pt}{4.818pt}}
\multiput(338.92,472.00)(-1.000,10.000){2}{\rule{1.000pt}{2.409pt}}
\put(339.42,508){\rule{1.000pt}{3.132pt}}
\multiput(337.92,508.00)(3.000,6.500){2}{\rule{1.000pt}{1.566pt}}
\put(342.92,521){\rule{1.000pt}{2.650pt}}
\multiput(340.92,521.00)(4.000,5.500){2}{\rule{1.000pt}{1.325pt}}
\multiput(348.86,532.00)(0.424,0.320){2}{\rule{0.102pt}{1.850pt}}
\multiput(344.92,532.00)(5.000,4.160){2}{\rule{1.000pt}{0.925pt}}
\multiput(353.84,540.00)(0.462,0.476){4}{\rule{0.111pt}{1.583pt}}
\multiput(349.92,540.00)(6.000,4.714){2}{\rule{1.000pt}{0.792pt}}
\multiput(358.00,549.84)(0.373,0.462){4}{\rule{1.417pt}{0.111pt}}
\multiput(358.00,545.92)(4.060,6.000){2}{\rule{0.708pt}{1.000pt}}
\multiput(365.00,555.84)(0.476,0.462){4}{\rule{1.583pt}{0.111pt}}
\multiput(365.00,551.92)(4.714,6.000){2}{\rule{0.792pt}{1.000pt}}
\put(373,559.92){\rule{2.168pt}{1.000pt}}
\multiput(373.00,557.92)(4.500,4.000){2}{\rule{1.084pt}{1.000pt}}
\multiput(382.00,565.86)(0.490,0.424){2}{\rule{2.050pt}{0.102pt}}
\multiput(382.00,561.92)(4.745,5.000){2}{\rule{1.025pt}{1.000pt}}
\put(391,568.92){\rule{2.168pt}{1.000pt}}
\multiput(391.00,566.92)(4.500,4.000){2}{\rule{1.084pt}{1.000pt}}
\put(400,572.92){\rule{2.650pt}{1.000pt}}
\multiput(400.00,570.92)(5.500,4.000){2}{\rule{1.325pt}{1.000pt}}
\put(411,576.92){\rule{2.409pt}{1.000pt}}
\multiput(411.00,574.92)(5.000,4.000){2}{\rule{1.204pt}{1.000pt}}
\put(421,580.92){\rule{2.891pt}{1.000pt}}
\multiput(421.00,578.92)(6.000,4.000){2}{\rule{1.445pt}{1.000pt}}
\put(433,584.92){\rule{2.650pt}{1.000pt}}
\multiput(433.00,582.92)(5.500,4.000){2}{\rule{1.325pt}{1.000pt}}
\put(444,588.92){\rule{2.891pt}{1.000pt}}
\multiput(444.00,586.92)(6.000,4.000){2}{\rule{1.445pt}{1.000pt}}
\put(456,592.92){\rule{3.132pt}{1.000pt}}
\multiput(456.00,590.92)(6.500,4.000){2}{\rule{1.566pt}{1.000pt}}
\multiput(469.00,598.86)(1.169,0.424){2}{\rule{2.850pt}{0.102pt}}
\multiput(469.00,594.92)(7.085,5.000){2}{\rule{1.425pt}{1.000pt}}
\multiput(482.00,603.86)(1.509,0.424){2}{\rule{3.250pt}{0.102pt}}
\multiput(482.00,599.92)(8.254,5.000){2}{\rule{1.625pt}{1.000pt}}
\multiput(497.00,608.86)(1.339,0.424){2}{\rule{3.050pt}{0.102pt}}
\multiput(497.00,604.92)(7.670,5.000){2}{\rule{1.525pt}{1.000pt}}
\multiput(511.00,613.84)(1.195,0.462){4}{\rule{2.750pt}{0.111pt}}
\multiput(511.00,609.92)(9.292,6.000){2}{\rule{1.375pt}{1.000pt}}
\multiput(526.00,619.84)(1.095,0.475){6}{\rule{2.536pt}{0.114pt}}
\multiput(526.00,615.92)(10.737,7.000){2}{\rule{1.268pt}{1.000pt}}
\multiput(542.00,626.84)(1.176,0.475){6}{\rule{2.679pt}{0.114pt}}
\multiput(542.00,622.92)(11.440,7.000){2}{\rule{1.339pt}{1.000pt}}
\multiput(559.00,633.84)(1.258,0.475){6}{\rule{2.821pt}{0.114pt}}
\multiput(559.00,629.92)(12.144,7.000){2}{\rule{1.411pt}{1.000pt}}
\multiput(577.00,640.83)(1.022,0.485){10}{\rule{2.361pt}{0.117pt}}
\multiput(577.00,636.92)(14.099,9.000){2}{\rule{1.181pt}{1.000pt}}
\multiput(596.00,649.83)(1.082,0.485){10}{\rule{2.472pt}{0.117pt}}
\multiput(596.00,645.92)(14.869,9.000){2}{\rule{1.236pt}{1.000pt}}
\multiput(616.00,658.83)(0.974,0.489){14}{\rule{2.250pt}{0.118pt}}
\multiput(616.00,654.92)(17.330,11.000){2}{\rule{1.125pt}{1.000pt}}
\multiput(638.00,669.83)(0.890,0.491){16}{\rule{2.083pt}{0.118pt}}
\multiput(638.00,665.92)(17.676,12.000){2}{\rule{1.042pt}{1.000pt}}
\multiput(660.00,681.83)(0.940,0.492){18}{\rule{2.173pt}{0.118pt}}
\multiput(660.00,677.92)(20.490,13.000){2}{\rule{1.087pt}{1.000pt}}
\multiput(685.00,694.83)(0.908,0.492){20}{\rule{2.107pt}{0.119pt}}
\multiput(685.00,690.92)(21.627,14.000){2}{\rule{1.054pt}{1.000pt}}
\multiput(711.00,708.83)(0.915,0.493){22}{\rule{2.117pt}{0.119pt}}
\multiput(711.00,704.92)(23.607,15.000){2}{\rule{1.058pt}{1.000pt}}
\multiput(739.00,723.83)(0.865,0.494){26}{\rule{2.015pt}{0.119pt}}
\multiput(739.00,719.92)(25.818,17.000){2}{\rule{1.007pt}{1.000pt}}
\multiput(769.00,740.83)(0.873,0.495){28}{\rule{2.028pt}{0.119pt}}
\multiput(769.00,736.92)(27.791,18.000){2}{\rule{1.014pt}{1.000pt}}
\multiput(801.00,758.83)(0.930,0.495){28}{\rule{2.139pt}{0.119pt}}
\multiput(801.00,754.92)(29.561,18.000){2}{\rule{1.069pt}{1.000pt}}
\multiput(835.00,776.83)(0.961,0.495){30}{\rule{2.197pt}{0.119pt}}
\multiput(835.00,772.92)(32.439,19.000){2}{\rule{1.099pt}{1.000pt}}
\multiput(872.00,795.83)(1.211,0.494){24}{\rule{2.688pt}{0.119pt}}
\multiput(872.00,791.92)(33.422,16.000){2}{\rule{1.344pt}{1.000pt}}
\multiput(911.00,811.83)(1.463,0.492){20}{\rule{3.179pt}{0.119pt}}
\multiput(911.00,807.92)(34.403,14.000){2}{\rule{1.589pt}{1.000pt}}
\multiput(952.00,825.83)(2.459,0.485){10}{\rule{5.028pt}{0.117pt}}
\multiput(952.00,821.92)(32.565,9.000){2}{\rule{2.514pt}{1.000pt}}
\put(995,832.42){\rule{10.359pt}{1.000pt}}
\multiput(995.00,830.92)(21.500,3.000){2}{\rule{5.179pt}{1.000pt}}
\multiput(1038.00,833.69)(4.071,-0.462){4}{\rule{7.417pt}{0.111pt}}
\multiput(1038.00,833.92)(27.606,-6.000){2}{\rule{3.708pt}{1.000pt}}
\multiput(1081.00,827.68)(1.307,-0.494){24}{\rule{2.875pt}{0.119pt}}
\multiput(1081.00,827.92)(36.033,-16.000){2}{\rule{1.438pt}{1.000pt}}
\multiput(1123.00,811.68)(0.800,-0.496){40}{\rule{1.875pt}{0.120pt}}
\multiput(1123.00,811.92)(35.108,-24.000){2}{\rule{0.937pt}{1.000pt}}
\multiput(1162.00,787.68)(0.534,-0.497){58}{\rule{1.341pt}{0.120pt}}
\multiput(1162.00,787.92)(33.217,-33.000){2}{\rule{0.670pt}{1.000pt}}
\multiput(1199.83,750.87)(0.497,-0.601){54}{\rule{0.120pt}{1.476pt}}
\multiput(1195.92,753.94)(31.000,-34.937){2}{\rule{1.000pt}{0.738pt}}
\multiput(1230.83,711.58)(0.497,-0.758){44}{\rule{0.120pt}{1.788pt}}
\multiput(1226.92,715.29)(26.000,-36.288){2}{\rule{1.000pt}{0.894pt}}
\multiput(1256.83,669.66)(0.495,-0.989){32}{\rule{0.119pt}{2.250pt}}
\multiput(1252.92,674.33)(20.000,-35.330){2}{\rule{1.000pt}{1.125pt}}
\multiput(1276.83,628.36)(0.494,-1.146){24}{\rule{0.119pt}{2.562pt}}
\multiput(1272.92,633.68)(16.000,-31.681){2}{\rule{1.000pt}{1.281pt}}
\multiput(1292.83,589.89)(0.491,-1.326){16}{\rule{0.118pt}{2.917pt}}
\multiput(1288.92,595.95)(12.000,-25.946){2}{\rule{1.000pt}{1.458pt}}
\multiput(1304.83,553.91)(0.481,-1.845){8}{\rule{0.116pt}{3.875pt}}
\multiput(1300.92,561.96)(8.000,-20.957){2}{\rule{1.000pt}{1.938pt}}
\put(1310.92,518){\rule{1.000pt}{5.541pt}}
\multiput(1308.92,529.50)(4.000,-11.500){2}{\rule{1.000pt}{2.770pt}}
\put(1313.42,498){\rule{1.000pt}{4.818pt}}
\multiput(1312.92,508.00)(1.000,-10.000){2}{\rule{1.000pt}{2.409pt}}
\put(340.0,492.0){\rule[-0.500pt]{1.000pt}{3.854pt}}
\put(1312.42,469){\rule{1.000pt}{3.132pt}}
\multiput(1313.92,475.50)(-3.000,-6.500){2}{\rule{1.000pt}{1.566pt}}
\put(1308.92,458){\rule{1.000pt}{2.650pt}}
\multiput(1310.92,463.50)(-4.000,-5.500){2}{\rule{1.000pt}{1.325pt}}
\multiput(1306.71,450.32)(-0.424,-0.320){2}{\rule{0.102pt}{1.850pt}}
\multiput(1306.92,454.16)(-5.000,-4.160){2}{\rule{1.000pt}{0.925pt}}
\multiput(1301.69,443.43)(-0.462,-0.476){4}{\rule{0.111pt}{1.583pt}}
\multiput(1301.92,446.71)(-6.000,-4.714){2}{\rule{1.000pt}{0.792pt}}
\multiput(1292.12,439.69)(-0.373,-0.462){4}{\rule{1.417pt}{0.111pt}}
\multiput(1295.06,439.92)(-4.060,-6.000){2}{\rule{0.708pt}{1.000pt}}
\put(1316.0,482.0){\rule[-0.500pt]{1.000pt}{3.854pt}}
\end{picture}


\vspace{1cm}
\caption{Solutions of van der Pol's equation for different values 
         of $\varepsilon^2$.}
\end{figure}

Besides a solution of equation we can create also the expressions
for the origin conditions which provide the periodic solution of
(\ref{eq-newvdp}) as a function (series) in $\varepsilon$. Really,
we have periodic solution when the expression $Z_1 Z_2$ is a constant
in time. Let us inverse series (\ref{n-trans}), (\ref{line-norm}) and 
change $Z_i$ in this production for their expressions via $\chi,d\chi/dt$
and $\varepsilon$. The result must be constant along the trajectory 
of the limit cycle too. Therefore we have equation of that trajectory
(there is no a clear dependence on time below, but $\chi$ and $d\chi/dt$
depend on time by implicit way):

$$
{d \over dt} [Z_1(\chi,d\chi/dt,\varepsilon) Z_2(\chi,d\chi/dt,\varepsilon)] 
= 0
$$
\noindent
 Differenting symbolically the above using the van der Pol's equation  
for an evaluation of the second derivation $d^2\chi/dt^2$ we will
have the corresponding equation in the form of the series:

$$
\sum_{i,j,k=0,1,\ldots} p_{i,j,k} \ \varepsilon^{2k} \chi^i
{d\chi \over dt}^j = 0, \ \ p_{0,0,0} \stackrel{\rm def}{=} 0
$$

 Suppose we are interested in the dependence of the origin condition
for the derivation in (\ref{eq-newvdp}) ${d\chi \over dt}|_{t=t_0}$ on 
$\varepsilon$ if $\chi|_{t=t_0} = 0$. In such a case we substitute zero 
for $\chi$ in the sum above and solve the equation resulted. The result for 
van der Pol's equation is:

\begin{equation}
\label{orcond-dChi}
\begin{array}{ll}

\chi = &0\\
{\mathstrut d\chi \over \mathstrut dt} = &2
+{\mathstrut 17 \over \mathstrut 96} \varepsilon^4
-{\mathstrut 1577 \over \mathstrut 552960} \varepsilon^8
-{\mathstrut 102956839 \over \mathstrut 55738368000} \varepsilon^{12}
+{\mathstrut 48722480822161 \over \mathstrut 157315969843200000}
\varepsilon^{16} + \ldots

\end{array}
\end{equation}

For the couple $\chi|_{t=t_0}, {d\chi \over dt}|_{t=t_0}=0$ we have:

\begin{equation}
\label{orcond-Chi}
\begin{array}{ll}

\chi = &2
+{\mathstrut 1 \over \mathstrut 96} \varepsilon^4
-{\mathstrut 1033 \over \mathstrut 552960} \varepsilon^8
+{\mathstrut 1019689 \over \mathstrut 55738368000} \varepsilon^{12}
+{\mathstrut 9835512276689 \over \mathstrut 157315969843200000} 
\varepsilon^{16} + \ldots\\
{\mathstrut d\chi \over \mathstrut dt} = &0
\end{array}
\end{equation}

\section{Henon--Heiles' equations as an example of the forth order 
         system}

   Henon--Heiles' system originally appeared in the theory of 
motion of particle in an axial symmetric gravitation field, for 
example in the problem of motion of stars in the Galactic field
(Henon, Heiles, 1964).
It is a system of two second order differential equations:


\begin{equation}
\label{eq-heiles}
\begin{array}{lcl}
\frac{d^{2}X}{\mathstrut dt^2}& = & - X - 2 X Y \\
\frac{\mathstrut d^{2}Y}{dt^2}& = & - Y - X^{2} + Y^{2}
\end{array}
\end{equation}

 
   System (\ref{eq-heiles}) can be rewritten in form 
(\ref{eq-general}) by the linear complex substitution:

$$
\begin{array}{ccc}
X & = & Y_1 + Y_2 \\
Y & = & Y_3 + Y_4 \\
\frac{\mathstrut dX}{\mathstrut dt} & = & i \ (Y_1 - Y_2) \\
\frac{\mathstrut dY}{\mathstrut dt} & = & i \ (Y_3 - Y_4)
\end{array}
$$

\noindent
after which we will have instead of (\ref{eq-heiles}):

\begin{equation}
\label{heiles-planar}
\begin{array}{lcrl}
\frac{\mathstrut dY_1}{\mathstrut dt}& = &  i &\!\!\!\! Y_1 + i \ (Y_1 + Y_2)\ 
                                                                 (Y_3 + Y_4)\\
\frac{\mathstrut dY_2}{\mathstrut dt}& = & -i &\!\!\!\! Y_2 - i \ (Y_1 + Y_2)\ 
                                                                 (Y_3 + Y_4)\\
\frac{\mathstrut dY_3}{\mathstrut dt}& = &  i &\!\!\!\! Y_3 + {i \over 2}\ [  
                                                                 (Y_1 + Y_2)^2 
                                                           - (Y_3 + Y_4)^2 ]\\
\frac{\mathstrut dY_4}{\mathstrut dt}& = & -i &\!\!\!\! Y_4 - {i \over 2}\ [  
                                                                 (Y_1 + Y_2)^2 
                                                           - (Y_3 + Y_4)^2 ]
\end{array}
\end{equation}

   The non-linear normalization for (\ref{eq-heiles}) is produced by the 
transformation:


\begin{equation}
\label{normal-tran}
Y_i = Z_i + Z_i
                       \!\!\!\!\!\!\!\!\!\!\!\!\!\!\!\!\!
                       \!\!\!\!\!\!\!\!\!\!\!\!\!\!\!\!\!
                       \sum_{\begin{array}{c}
                                    p_i \geq -1 \\
                           p_1,\ldots,p_{i-1},p_{i+1},\ldots,p_4 \geq 0 
                                          \end{array}}
                                   \!\!\!\!\!\!\!\!\!\!\!\!\!\!\!\!
                                   \!\!\!\!\!\!\!\!\!\!\!\!\!\!\!\!
                                    h_{i,p_1,p_2,p_3,p_4}
                           \ Z_1^{p_1}Z_2^{p_2}Z_3^{p_3}Z_4^{p_4},
\quad i=1,\ldots ,4
\end{equation}

    The vector of eigenvalues of the linear part will consist of two
 couples of
complex conjugated imaginary units: $ \overline{\lambda} = 
\{i,-i,i,-i\}$. I.e.
the normal form for (\ref{heiles-planar}) will be:


\begin{equation}
\label{heiles-normal}
\frac{dZ_i}{dt} = \lambda_iZ_i + Z_i
                                   \!\!\!\!\!\!\!\!\!\!\!\!\!\!\!\!\!
                                   \!\!\!\!\!\!\!\!\!\!\!\!\!\!\!\!\!
                                    \sum_{\begin{array}{c}
                                    p_i \geq -1 \\
                           p_1,\ldots,p_{i-1},p_{i+1},\ldots,p_4 \geq 0 \\
                                       p_1 + p_3 = p_2 + p_4
                                          \end{array}}
                                   \!\!\!\!\!\!\!\!\!\!\!\!\!\!\!\!
                                   \!\!\!\!\!\!\!\!\!\!\!\!\!\!\!\!
                                    g_{i,p_1,p_2,p_3,p_4}
                           \ Z_1^{p_1}Z_2^{p_2}Z_3^{p_3}Z_4^{p_4},
\quad i=1,\ldots ,4
\end{equation}

 \subsection{periodic solutions of Henon--Heiles' system}

   The domain in phase space, where the normalizing transformation  will
be analytic for a system which has two couples of complex conjugated 
imaginary units as eigenvalues, is defined by equations 
(\ref{cond-A}):

\begin{equation}
\label{heiles-condition}
\lambda_iZ_i a = \lambda_iZ_i + Z_i
                                   \!\!\!\!\!\!\!\!\!\!\!\!\!\!\!\!\!
                                   \!\!\!\!\!\!\!\!\!\!\!\!\!\!\!\!\!
                                   \sum_{\begin{array}{c}
                                    p_i \geq -1 \\
                           p_1,\ldots,p_{i-1},p_{i+1},\ldots,p_4 \geq 0 \\
                                       p_1 + p_3 = p_2 + p_4
                                          \end{array}}
                                   \!\!\!\!\!\!\!\!\!\!\!\!\!\!\!\!
                                   \!\!\!\!\!\!\!\!\!\!\!\!\!\!\!\!
                                    g_{i,p_1,p_2,p_3,p_4}
                           \ Z_1^{p_1}Z_2^{p_2}Z_3^{p_3}Z_4^{p_4},
\quad i=1,\ldots ,4
\end{equation}

\noindent
where $a$ is an expression, which does not depend on $i$ index. It
can be proved that $a$ does not depend on time if (\ref{heiles-condition})
are satisfied.

   The set in phase space of (\ref{heiles-normal}) for which 
(\ref{heiles-condition}) is fulfilled will produce by the inverse
transformation to 
(\ref{normal-tran}) the invariant periodic set of (\ref{heiles-planar}) 
that contracts to the stationary point (a zero point here) (Bruno, 1993).

   Equations (\ref{heiles-condition}) are simplified using Groebner
bases technique (Buchberger, 1985), (Boege, Gebauer, Kredel, 1986) 
with the REDUCE (Hearn, 1987) computer algebra system. 
We obtained 8 proper solutions of (\ref{heiles-condition})
and they produce 6 families of different periodic solutions of 
(\ref{eq-heiles}) with 4 different values of a frequency.

As 4 equations (\ref{heiles-condition}) are not independent we have as 
a solution a hyperplane rather than a point:

$$
\begin{array}{lll}
\mbox{solution 1: } &  Z_1 Z_4 = \ \ Z_2 Z_3, & Z_3^2 = 3 Z_1^2; \\
\mbox{solution 2: } &  Z_1 Z_4 = \ \ Z_2 Z_3, & Z_1^2 = 3 Z_3^2; \\
\mbox{solution 3: } &  Z_1 Z_4 = - Z_2 Z_3, & Z_1^2 = \ \ \! Z_3^2;\\
\mbox{solution 4: } &  Z_3 = Z_4 = 0;& \\
\mbox{solution 5: } &  Z_1 = Z_2 = 0;& \\
\mbox{solution 6: } &  Z_2 = Z_4 = 0;& \\
\mbox{solution 7: } &  Z_1 = Z_3 = 0;& 
\end{array}
$$

   Corresponding frequences for the solutions of (\ref{eq-heiles}) are:
$$
\begin{array}{l}

% case 3: z3:=z4:=0, case 6: z1*z4:=z2*z3, z3^2:=3*z1^2; sol 4,1:
\omega_{1,4} = 
1 -{\mathstrut 5 \over \mathstrut 6} H + 
{17 \over 48} H^2 + {127517 \over 38880} H^3 + 
{\mathstrut 51952319 \over \mathstrut 3732480} H^4 +
{1675438657 \over 111974400} H^5 \\

% case 5: z1:=z2:=0, case 7: z1*z4:=z2*z3, z1^2:=3*z3^2; sol 5,2:
\omega_{2,5} = 
1 - {\mathstrut 5 \over \mathstrut 6} H - 
{95 \over 48} H^2 - {54935 \over 7776} H^3 - 
{\mathstrut 22030445 \over \mathstrut 746496} H^4 - 
{200207485 \over 1492992} H^5 \\

% case 1: z1*z4:=-z2*z3, z1^2:=z3^2; sol 3:
\omega_{3} = 
1 + {\mathstrut 1 \over \mathstrut 3} H - 
{2 \over 3} H^2 + {5389 \over 2430} H^3 - 
{\mathstrut 52393 \over \mathstrut 5832} H4 + 
{29471957 \over 729000} H^5 \\

% case 2: z2:=z4:=0 and case 4: z1:=z3:=0; sol 6,7: 
\omega_{6} = \omega_{7} = 1 

\end{array}
$$

\noindent
where $H$ is the full energy:

$$
H = {1 \over 2}[(\frac{dX}{dt})^2 + (\frac{dY}{dt})^2
                + X^2 + Y^2] + X^2 Y - {1 \over 3} Y^3
$$
  

   For each frequence the above mentioned the corresponding solution 
in analytical
form (as a Fourier series in $t$ and $H$) was obtained. This is the
solution which corresponds to $\omega_1$:

$$
\begin{array}{rl}
X = \hfill &\\
\frac{\sqrt{3}}{\mathstrut 2799360}\cos(8 \omega t)\varepsilon^8&  + \\
\frac{\mathstrut 71\sqrt{2}}{\mathstrut 18662400}\cos(7 \omega t)\varepsilon^7& 
+ \\
\frac{\mathstrut \sqrt{3}}{\mathstrut 259200}\cos(6 \omega t)\varepsilon^6
&{}\!\!\!\!\!(130\varepsilon^2 + 9)  + \\
\frac{\mathstrut 17\sqrt{2}}{\mathstrut 1555200}\cos(5 \omega t)\varepsilon^5
&{}\!\!\!\!\!(341\varepsilon^2 + 30)  +\\ 
\frac{\mathstrut \sqrt{3}}{\mathstrut 583200}\cos(4 \omega t)\varepsilon^4
&{}\!\!\!\!\!(3451\varepsilon^4 + 3204\varepsilon^2 + 540) + \\
\frac{\mathstrut \sqrt{2}}{\mathstrut 3110400 }\cos(3 \omega t)\varepsilon^3
&{}\!\!\!\!\!(1238051\varepsilon^4 + 398160\varepsilon^2 + 64800) + \\
\frac{\mathstrut \sqrt{3}}{1399680}\cos(2 \omega t)\varepsilon^2
&{}\!\!\!\!\!(68930\varepsilon^6 
            + 2152809\varepsilon^4 \\
&           + 933120\varepsilon^2  + 233280)  + \\
\frac{\sqrt{2}}{18662400}\cos( \omega t)\varepsilon
&{}\!\!\!\!\!(- 51700693\varepsilon^6
            - 5862600\varepsilon^4 \\
&           + 4795200\varepsilon^2 + 9331200)  + \\
\frac{\sqrt{3}}{1555200}\varepsilon^2
&{}\!\!\!\!\!(8445851\varepsilon^6 + 1072800\varepsilon^4 \\ 
&           + 21600\varepsilon^2 - 777600)\\
&\\
Y =\hfill&- \\
\frac{1}{\mathstrut 2799360}\cos(8 \omega t)\varepsilon^8&  + \\
\frac{\mathstrut 71\sqrt{6}}{\mathstrut 18662400}\cos(7 \omega t)\varepsilon^7& 
- \\
\frac{\mathstrut 1}{\mathstrut 259200}\cos(6 \omega t)\varepsilon^6
&{}\!\!\!\!\!(130\varepsilon^2 + 9)  + \\ 
\frac{\mathstrut 17\sqrt{6}}{\mathstrut 1555200}\cos(5 \omega t)\varepsilon^5
&{}\!\!\!\!\!(341\varepsilon^2 + 30) - \\
\frac{\mathstrut 1}{\mathstrut 583200}\cos(4 \omega t)\varepsilon^4
&{}\!\!\!\!\!(3451\varepsilon^4 + 3204\varepsilon^2 + 540)  + \\
\frac{\mathstrut \sqrt{6}}{\mathstrut 3110400}\cos(3 \omega t)\varepsilon^3
&{}\!\!\!\!\!(1238051\varepsilon^4 + 398160\varepsilon^2 + 64800) - \\
\frac{\mathstrut 1}{1399680}\cos(2 \omega t)\varepsilon^2
&{}\!\!\!\!\!(68930\varepsilon^6 + 2152809\varepsilon^4 \\ 
&       + 933120\varepsilon^2 + 233280)  - \\
\frac{\sqrt{6}}{18662400}\cos( \omega t)\varepsilon
&{}\!\!\!\!\!(51700693\varepsilon^6 
        + 5862600\varepsilon^4 \\
&       - 4795200\varepsilon^2 - 9331200) - \\
\frac{1}{1555200}\varepsilon^2
&{}\!\!\!\!\!(8445851\varepsilon^6 + 1072800\varepsilon^4 \\ 
&       + 21600\varepsilon^2 - 777600)
\end{array}
$$


Results  of the calculation were verified  by two ways. The first one 
was a direct substitution of the series in the original equations 
(\ref{eq-heiles}). All results of this procedure contained only terms 
with negligible orders. The second way was a comparison of the numerical 
solutions of equations (\ref{eq-heiles}) calculated by the Runge--Kutta
method with the  tabulated  values of  the approximate solutions at values
of energy $H = {1 \over 24}$ and $H = {1 \over 12}$. 
We evaluated the values of $f_{err}$ which are  maximum values of 
relative deviations of the numerical solutions from the 
approximate solutions in a phase space during one period. The results 
are:

$$
\begin{array}{lll}
      &H = \frac{1}{\mathstrut 24}: \qquad & H = \frac{1}{\mathstrut 12}: \\
\mbox{for solutions with }\omega_1 :
      \quad & f_{err} = 1.3\times 10^{-5}& f_{err} = 7.6\times 10^{-4}\\
\mbox{for solutions with }\omega_2 :
      \quad & f_{err} = 4.7\times 10^{-5}& f_{err} = 4.4\times 10^{-3}\\
\mbox{for solutions with }\omega_3 :
      \quad & f_{err} = 1.1\times 10^{-5}& f_{err} = 5.0\times 10^{-4}\\
\mbox{for solutions with }\omega_4 :
      \quad & f_{err} = 1.3\times 10^{-5}& f_{err} = 7.6\times 10^{-4}\\
\mbox{for solutions with }\omega_5 :
      \quad & f_{err} = 2.9\times 10^{-5}& f_{err} = 1.8\times 10^{-3}

\end{array}
$$

\noindent
We did not check numerically the complex solutions with $\omega_6,
\omega_7$.

  In Fig.\ref{fig-henon} intersections of periodic solutions of 
Henon--Heiles' system with plane $X = 0,  
\frac{dX}{dt}=\frac{dX}{dt}(X=0,Y,Y',H)$ for the full energy 
${1 \over 12}$ are displayed. Solution 5 lies in that plane and
other solutions cross this plane in a couple points each except
solution 4 which looks like a digit "8" and crosses the plane $X=0$ 
in a single point.

\begin{figure}
\label{fig-henon}
\vspace{1cm}
\fbox{$\frac{d^{2}X}{\mathstrut dt^2} =  - X - 2 X Y, \ \ 
\frac{\mathstrut d^{2}Y}{dt^2} =  - Y - X^{2} + Y^{2}$}

% GNUPLOT: LaTeX picture
\setlength{\unitlength}{0.240900pt}
\ifx\plotpoint\undefined\newsavebox{\plotpoint}\fi
\sbox{\plotpoint}{\rule[-0.500pt]{1.000pt}{1.000pt}}%
\begin{picture}(1500,900)(0,0)
\font\gnuplot=cmr10 at 10pt
\gnuplot
\sbox{\plotpoint}{\rule[-0.500pt]{1.000pt}{1.000pt}}%
\put(220.0,495.0){\rule[-0.500pt]{292.934pt}{1.000pt}}
\put(760.0,113.0){\rule[-0.500pt]{1.000pt}{184.048pt}}
\put(220.0,113.0){\rule[-0.500pt]{4.818pt}{1.000pt}}
\put(198,113){\makebox(0,0)[r]{-0.5}}
\put(1416.0,113.0){\rule[-0.500pt]{4.818pt}{1.000pt}}
\put(220.0,189.0){\rule[-0.500pt]{4.818pt}{1.000pt}}
\put(198,189){\makebox(0,0)[r]{-0.4}}
\put(1416.0,189.0){\rule[-0.500pt]{4.818pt}{1.000pt}}
\put(220.0,266.0){\rule[-0.500pt]{4.818pt}{1.000pt}}
\put(198,266){\makebox(0,0)[r]{-0.3}}
\put(1416.0,266.0){\rule[-0.500pt]{4.818pt}{1.000pt}}
\put(220.0,342.0){\rule[-0.500pt]{4.818pt}{1.000pt}}
\put(198,342){\makebox(0,0)[r]{-0.2}}
\put(1416.0,342.0){\rule[-0.500pt]{4.818pt}{1.000pt}}
\put(220.0,419.0){\rule[-0.500pt]{4.818pt}{1.000pt}}
\put(198,419){\makebox(0,0)[r]{-0.1}}
\put(1416.0,419.0){\rule[-0.500pt]{4.818pt}{1.000pt}}
\put(220.0,495.0){\rule[-0.500pt]{4.818pt}{1.000pt}}
\put(198,495){\makebox(0,0)[r]{0}}
\put(1416.0,495.0){\rule[-0.500pt]{4.818pt}{1.000pt}}
\put(220.0,571.0){\rule[-0.500pt]{4.818pt}{1.000pt}}
\put(198,571){\makebox(0,0)[r]{0.1}}
\put(1416.0,571.0){\rule[-0.500pt]{4.818pt}{1.000pt}}
\put(220.0,648.0){\rule[-0.500pt]{4.818pt}{1.000pt}}
\put(198,648){\makebox(0,0)[r]{0.2}}
\put(1416.0,648.0){\rule[-0.500pt]{4.818pt}{1.000pt}}
\put(220.0,724.0){\rule[-0.500pt]{4.818pt}{1.000pt}}
\put(198,724){\makebox(0,0)[r]{0.3}}
\put(1416.0,724.0){\rule[-0.500pt]{4.818pt}{1.000pt}}
\put(220.0,801.0){\rule[-0.500pt]{4.818pt}{1.000pt}}
\put(198,801){\makebox(0,0)[r]{0.4}}
\put(1416.0,801.0){\rule[-0.500pt]{4.818pt}{1.000pt}}
\put(220.0,877.0){\rule[-0.500pt]{4.818pt}{1.000pt}}
\put(198,877){\makebox(0,0)[r]{0.5}}
\put(1416.0,877.0){\rule[-0.500pt]{4.818pt}{1.000pt}}
\put(220.0,113.0){\rule[-0.500pt]{1.000pt}{4.818pt}}
\put(220,68){\makebox(0,0){-0.4}}
\put(220.0,857.0){\rule[-0.500pt]{1.000pt}{4.818pt}}
\put(355.0,113.0){\rule[-0.500pt]{1.000pt}{4.818pt}}
\put(355,68){\makebox(0,0){-0.3}}
\put(355.0,857.0){\rule[-0.500pt]{1.000pt}{4.818pt}}
\put(490.0,113.0){\rule[-0.500pt]{1.000pt}{4.818pt}}
\put(490,68){\makebox(0,0){-0.2}}
\put(490.0,857.0){\rule[-0.500pt]{1.000pt}{4.818pt}}
\put(625.0,113.0){\rule[-0.500pt]{1.000pt}{4.818pt}}
\put(625,68){\makebox(0,0){-0.1}}
\put(625.0,857.0){\rule[-0.500pt]{1.000pt}{4.818pt}}
\put(760.0,113.0){\rule[-0.500pt]{1.000pt}{4.818pt}}
\put(760,68){\makebox(0,0){0}}
\put(760.0,857.0){\rule[-0.500pt]{1.000pt}{4.818pt}}
\put(896.0,113.0){\rule[-0.500pt]{1.000pt}{4.818pt}}
\put(896,68){\makebox(0,0){0.1}}
\put(896.0,857.0){\rule[-0.500pt]{1.000pt}{4.818pt}}
\put(1031.0,113.0){\rule[-0.500pt]{1.000pt}{4.818pt}}
\put(1031,68){\makebox(0,0){0.2}}
\put(1031.0,857.0){\rule[-0.500pt]{1.000pt}{4.818pt}}
\put(1166.0,113.0){\rule[-0.500pt]{1.000pt}{4.818pt}}
\put(1166,68){\makebox(0,0){0.3}}
\put(1166.0,857.0){\rule[-0.500pt]{1.000pt}{4.818pt}}
\put(1301.0,113.0){\rule[-0.500pt]{1.000pt}{4.818pt}}
\put(1301,68){\makebox(0,0){0.4}}
\put(1301.0,857.0){\rule[-0.500pt]{1.000pt}{4.818pt}}
\put(1436.0,113.0){\rule[-0.500pt]{1.000pt}{4.818pt}}
\put(1436,68){\makebox(0,0){0.5}}
\put(1436.0,857.0){\rule[-0.500pt]{1.000pt}{4.818pt}}
\put(220.0,113.0){\rule[-0.500pt]{292.934pt}{1.000pt}}
\put(1436.0,113.0){\rule[-0.500pt]{1.000pt}{184.048pt}}
\put(220.0,877.0){\rule[-0.500pt]{292.934pt}{1.000pt}}
\put(111,495){\makebox(0,0){$\dot{Y}$}}
\put(762,23){\makebox(0,0){$Y$}}
\put(220.0,113.0){\rule[-0.500pt]{1.000pt}{184.048pt}}
\put(1306,812){\makebox(0,0)[r]{$\omega_1$}}
\put(1350,812){\makebox(0,0){$+$}}
\put(1040,281){\makebox(0,0){$+$}}
\put(1040,709){\makebox(0,0){$+$}}
\sbox{\plotpoint}{\rule[-0.175pt]{0.350pt}{0.350pt}}%
\put(1306,767){\makebox(0,0)[r]{$\omega_2$}}
\put(1350,767){\raisebox{-.8pt}{\makebox(0,0){$\Diamond$}}}
\put(759,339){\raisebox{-.8pt}{\makebox(0,0){$\Diamond$}}}
\put(759,651){\raisebox{-.8pt}{\makebox(0,0){$\Diamond$}}}
\sbox{\plotpoint}{\rule[-0.300pt]{0.600pt}{0.600pt}}%
\put(1306,722){\makebox(0,0)[r]{$\omega_3$}}
\put(1350,722){\raisebox{-.8pt}{\makebox(0,0){$\Box$}}}
\put(350,495){\raisebox{-.8pt}{\makebox(0,0){$\Box$}}}
\put(1105,495){\raisebox{-.8pt}{\makebox(0,0){$\Box$}}}
\sbox{\plotpoint}{\rule[-0.250pt]{0.500pt}{0.500pt}}%
\put(1306,677){\makebox(0,0)[r]{$\omega_4$}}
\put(1350,677){\makebox(0,0){$\times$}}
\put(598,495){\makebox(0,0){$\times$}}
\sbox{\plotpoint}{\rule[-0.500pt]{1.000pt}{1.000pt}}%
\put(1306,632){\makebox(0,0)[r]{$\omega_5$}}
\put(1328.0,632.0){\rule[-0.500pt]{15.899pt}{1.000pt}}
\put(1436,495){\usebox{\plotpoint}}
\put(1433.42,482){\rule{1.000pt}{3.132pt}}
\multiput(1433.92,488.50)(-1.000,-6.500){2}{\rule{1.000pt}{1.566pt}}
\put(1431.42,469){\rule{1.000pt}{3.132pt}}
\multiput(1432.92,475.50)(-3.000,-6.500){2}{\rule{1.000pt}{1.566pt}}
\put(1427.92,455){\rule{1.000pt}{3.373pt}}
\multiput(1429.92,462.00)(-4.000,-7.000){2}{\rule{1.000pt}{1.686pt}}
\multiput(1425.71,443.17)(-0.424,-1.169){2}{\rule{0.102pt}{2.850pt}}
\multiput(1425.92,449.08)(-5.000,-7.085){2}{\rule{1.000pt}{1.425pt}}
\multiput(1420.68,434.22)(-0.481,-0.745){8}{\rule{0.116pt}{1.875pt}}
\multiput(1420.92,438.11)(-8.000,-9.108){2}{\rule{1.000pt}{0.938pt}}
\multiput(1412.68,421.22)(-0.481,-0.745){8}{\rule{0.116pt}{1.875pt}}
\multiput(1412.92,425.11)(-8.000,-9.108){2}{\rule{1.000pt}{0.938pt}}
\multiput(1404.68,410.06)(-0.489,-0.543){14}{\rule{0.118pt}{1.432pt}}
\multiput(1404.92,413.03)(-11.000,-10.028){2}{\rule{1.000pt}{0.716pt}}
\multiput(1393.68,397.12)(-0.491,-0.541){16}{\rule{0.118pt}{1.417pt}}
\multiput(1393.92,400.06)(-12.000,-11.060){2}{\rule{1.000pt}{0.708pt}}
\multiput(1378.49,386.68)(-0.500,-0.492){18}{\rule{1.327pt}{0.118pt}}
\multiput(1381.25,386.92)(-11.246,-13.000){2}{\rule{0.663pt}{1.000pt}}
\multiput(1364.17,373.68)(-0.540,-0.492){18}{\rule{1.404pt}{0.118pt}}
\multiput(1367.09,373.92)(-12.086,-13.000){2}{\rule{0.702pt}{1.000pt}}
\multiput(1348.53,360.68)(-0.620,-0.492){18}{\rule{1.558pt}{0.118pt}}
\multiput(1351.77,360.92)(-13.767,-13.000){2}{\rule{0.779pt}{1.000pt}}
\multiput(1330.74,347.68)(-0.716,-0.491){16}{\rule{1.750pt}{0.118pt}}
\multiput(1334.37,347.92)(-14.368,-12.000){2}{\rule{0.875pt}{1.000pt}}
\multiput(1312.58,335.68)(-0.740,-0.492){18}{\rule{1.788pt}{0.118pt}}
\multiput(1316.29,335.92)(-16.288,-13.000){2}{\rule{0.894pt}{1.000pt}}
\multiput(1291.94,322.68)(-0.820,-0.492){18}{\rule{1.942pt}{0.118pt}}
\multiput(1295.97,322.92)(-17.969,-13.000){2}{\rule{0.971pt}{1.000pt}}
\multiput(1269.01,309.68)(-0.934,-0.491){16}{\rule{2.167pt}{0.118pt}}
\multiput(1273.50,309.92)(-18.503,-12.000){2}{\rule{1.083pt}{1.000pt}}
\multiput(1245.66,297.68)(-0.977,-0.491){16}{\rule{2.250pt}{0.118pt}}
\multiput(1250.33,297.92)(-19.330,-12.000){2}{\rule{1.125pt}{1.000pt}}
\multiput(1220.97,285.68)(-1.065,-0.491){16}{\rule{2.417pt}{0.118pt}}
\multiput(1225.98,285.92)(-20.984,-12.000){2}{\rule{1.208pt}{1.000pt}}
\multiput(1193.40,273.68)(-1.261,-0.489){14}{\rule{2.795pt}{0.118pt}}
\multiput(1199.20,273.92)(-22.198,-11.000){2}{\rule{1.398pt}{1.000pt}}
\multiput(1165.93,262.68)(-1.195,-0.491){16}{\rule{2.667pt}{0.118pt}}
\multiput(1171.47,262.92)(-23.465,-12.000){2}{\rule{1.333pt}{1.000pt}}
\multiput(1134.51,250.68)(-1.501,-0.487){12}{\rule{3.250pt}{0.117pt}}
\multiput(1141.25,250.92)(-23.254,-10.000){2}{\rule{1.625pt}{1.000pt}}
\multiput(1104.09,240.68)(-1.554,-0.487){12}{\rule{3.350pt}{0.117pt}}
\multiput(1111.05,240.92)(-24.047,-10.000){2}{\rule{1.675pt}{1.000pt}}
\multiput(1070.74,230.68)(-1.860,-0.485){10}{\rule{3.917pt}{0.117pt}}
\multiput(1078.87,230.92)(-24.871,-9.000){2}{\rule{1.958pt}{1.000pt}}
\multiput(1037.74,221.68)(-1.860,-0.485){10}{\rule{3.917pt}{0.117pt}}
\multiput(1045.87,221.92)(-24.871,-9.000){2}{\rule{1.958pt}{1.000pt}}
\multiput(1001.80,212.68)(-2.257,-0.481){8}{\rule{4.625pt}{0.116pt}}
\multiput(1011.40,212.92)(-25.401,-8.000){2}{\rule{2.312pt}{1.000pt}}
\multiput(964.21,204.69)(-2.641,-0.475){6}{\rule{5.250pt}{0.114pt}}
\multiput(975.10,204.92)(-24.103,-7.000){2}{\rule{2.625pt}{1.000pt}}
\multiput(924.36,197.69)(-3.455,-0.462){4}{\rule{6.417pt}{0.111pt}}
\multiput(937.68,197.92)(-23.682,-6.000){2}{\rule{3.208pt}{1.000pt}}
\put(877,189.92){\rule{8.913pt}{1.000pt}}
\multiput(895.50,191.92)(-18.500,-4.000){2}{\rule{4.457pt}{1.000pt}}
\put(840,185.92){\rule{8.913pt}{1.000pt}}
\multiput(858.50,187.92)(-18.500,-4.000){2}{\rule{4.457pt}{1.000pt}}
\put(802,182.92){\rule{9.154pt}{1.000pt}}
\multiput(821.00,183.92)(-19.000,-2.000){2}{\rule{4.577pt}{1.000pt}}
\put(764,181.42){\rule{9.154pt}{1.000pt}}
\multiput(783.00,181.92)(-19.000,-1.000){2}{\rule{4.577pt}{1.000pt}}
\put(726,181.42){\rule{9.154pt}{1.000pt}}
\multiput(745.00,180.92)(-19.000,1.000){2}{\rule{4.577pt}{1.000pt}}
\put(688,182.92){\rule{9.154pt}{1.000pt}}
\multiput(707.00,181.92)(-19.000,2.000){2}{\rule{4.577pt}{1.000pt}}
\put(651,185.92){\rule{8.913pt}{1.000pt}}
\multiput(669.50,183.92)(-18.500,4.000){2}{\rule{4.457pt}{1.000pt}}
\multiput(619.24,191.86)(-5.244,0.424){2}{\rule{7.650pt}{0.102pt}}
\multiput(635.12,187.92)(-21.122,5.000){2}{\rule{3.825pt}{1.000pt}}
\multiput(591.61,196.84)(-2.723,0.475){6}{\rule{5.393pt}{0.114pt}}
\multiput(602.81,192.92)(-24.807,7.000){2}{\rule{2.696pt}{1.000pt}}
\multiput(560.36,203.83)(-2.040,0.485){10}{\rule{4.250pt}{0.117pt}}
\multiput(569.18,199.92)(-27.179,9.000){2}{\rule{2.125pt}{1.000pt}}
\multiput(528.13,212.83)(-1.549,0.489){14}{\rule{3.341pt}{0.118pt}}
\multiput(535.07,208.92)(-27.066,11.000){2}{\rule{1.670pt}{1.000pt}}
\multiput(496.74,223.83)(-1.220,0.492){18}{\rule{2.712pt}{0.118pt}}
\multiput(502.37,219.92)(-26.372,13.000){2}{\rule{1.356pt}{1.000pt}}
\multiput(465.77,236.83)(-1.093,0.492){20}{\rule{2.464pt}{0.119pt}}
\multiput(470.89,232.92)(-25.885,14.000){2}{\rule{1.232pt}{1.000pt}}
\multiput(436.44,250.83)(-0.889,0.494){24}{\rule{2.063pt}{0.119pt}}
\multiput(440.72,246.92)(-24.719,16.000){2}{\rule{1.031pt}{1.000pt}}
\multiput(408.74,266.83)(-0.731,0.495){28}{\rule{1.750pt}{0.119pt}}
\multiput(412.37,262.92)(-23.368,18.000){2}{\rule{0.875pt}{1.000pt}}
\multiput(382.50,284.83)(-0.638,0.495){30}{\rule{1.566pt}{0.119pt}}
\multiput(385.75,280.92)(-21.750,19.000){2}{\rule{0.783pt}{1.000pt}}
\multiput(358.61,303.83)(-0.503,0.496){34}{\rule{1.298pt}{0.119pt}}
\multiput(361.31,299.92)(-19.307,21.000){2}{\rule{0.649pt}{1.000pt}}
\multiput(339.68,323.00)(-0.495,0.529){32}{\rule{0.119pt}{1.350pt}}
\multiput(339.92,323.00)(-20.000,19.198){2}{\rule{1.000pt}{0.675pt}}
\multiput(319.68,345.00)(-0.494,0.654){26}{\rule{0.119pt}{1.603pt}}
\multiput(319.92,345.00)(-17.000,19.673){2}{\rule{1.000pt}{0.801pt}}
\multiput(302.68,368.00)(-0.492,0.834){20}{\rule{0.119pt}{1.964pt}}
\multiput(302.92,368.00)(-14.000,19.923){2}{\rule{1.000pt}{0.982pt}}
\multiput(288.68,392.00)(-0.489,1.118){14}{\rule{0.118pt}{2.523pt}}
\multiput(288.92,392.00)(-11.000,19.764){2}{\rule{1.000pt}{1.261pt}}
\multiput(277.68,417.00)(-0.481,1.639){8}{\rule{0.116pt}{3.500pt}}
\multiput(277.92,417.00)(-8.000,18.736){2}{\rule{1.000pt}{1.750pt}}
\multiput(269.71,443.00)(-0.424,3.377){2}{\rule{0.102pt}{5.450pt}}
\multiput(269.92,443.00)(-5.000,14.688){2}{\rule{1.000pt}{2.725pt}}
\put(264.42,469){\rule{1.000pt}{6.263pt}}
\multiput(264.92,469.00)(-1.000,13.000){2}{\rule{1.000pt}{3.132pt}}
\put(264.42,495){\rule{1.000pt}{6.263pt}}
\multiput(263.92,495.00)(1.000,13.000){2}{\rule{1.000pt}{3.132pt}}
\multiput(268.86,521.00)(0.424,3.377){2}{\rule{0.102pt}{5.450pt}}
\multiput(264.92,521.00)(5.000,14.688){2}{\rule{1.000pt}{2.725pt}}
\multiput(273.83,547.00)(0.481,1.639){8}{\rule{0.116pt}{3.500pt}}
\multiput(269.92,547.00)(8.000,18.736){2}{\rule{1.000pt}{1.750pt}}
\multiput(281.83,573.00)(0.489,1.118){14}{\rule{0.118pt}{2.523pt}}
\multiput(277.92,573.00)(11.000,19.764){2}{\rule{1.000pt}{1.261pt}}
\multiput(292.83,598.00)(0.492,0.834){20}{\rule{0.119pt}{1.964pt}}
\multiput(288.92,598.00)(14.000,19.923){2}{\rule{1.000pt}{0.982pt}}
\multiput(306.83,622.00)(0.494,0.654){26}{\rule{0.119pt}{1.603pt}}
\multiput(302.92,622.00)(17.000,19.673){2}{\rule{1.000pt}{0.801pt}}
\multiput(323.83,645.00)(0.495,0.529){32}{\rule{0.119pt}{1.350pt}}
\multiput(319.92,645.00)(20.000,19.198){2}{\rule{1.000pt}{0.675pt}}
\multiput(342.00,668.83)(0.503,0.496){34}{\rule{1.298pt}{0.119pt}}
\multiput(342.00,664.92)(19.307,21.000){2}{\rule{0.649pt}{1.000pt}}
\multiput(364.00,689.83)(0.638,0.495){30}{\rule{1.566pt}{0.119pt}}
\multiput(364.00,685.92)(21.750,19.000){2}{\rule{0.783pt}{1.000pt}}
\multiput(389.00,708.83)(0.731,0.495){28}{\rule{1.750pt}{0.119pt}}
\multiput(389.00,704.92)(23.368,18.000){2}{\rule{0.875pt}{1.000pt}}
\multiput(416.00,726.83)(0.889,0.494){24}{\rule{2.063pt}{0.119pt}}
\multiput(416.00,722.92)(24.719,16.000){2}{\rule{1.031pt}{1.000pt}}
\multiput(445.00,742.83)(1.093,0.492){20}{\rule{2.464pt}{0.119pt}}
\multiput(445.00,738.92)(25.885,14.000){2}{\rule{1.232pt}{1.000pt}}
\multiput(476.00,756.83)(1.220,0.492){18}{\rule{2.712pt}{0.118pt}}
\multiput(476.00,752.92)(26.372,13.000){2}{\rule{1.356pt}{1.000pt}}
\multiput(508.00,769.83)(1.549,0.489){14}{\rule{3.341pt}{0.118pt}}
\multiput(508.00,765.92)(27.066,11.000){2}{\rule{1.670pt}{1.000pt}}
\multiput(542.00,780.83)(2.040,0.485){10}{\rule{4.250pt}{0.117pt}}
\multiput(542.00,776.92)(27.179,9.000){2}{\rule{2.125pt}{1.000pt}}
\multiput(578.00,789.84)(2.723,0.475){6}{\rule{5.393pt}{0.114pt}}
\multiput(578.00,785.92)(24.807,7.000){2}{\rule{2.696pt}{1.000pt}}
\multiput(614.00,796.86)(5.244,0.424){2}{\rule{7.650pt}{0.102pt}}
\multiput(614.00,792.92)(21.122,5.000){2}{\rule{3.825pt}{1.000pt}}
\put(651,799.92){\rule{8.913pt}{1.000pt}}
\multiput(651.00,797.92)(18.500,4.000){2}{\rule{4.457pt}{1.000pt}}
\put(688,802.92){\rule{9.154pt}{1.000pt}}
\multiput(688.00,801.92)(19.000,2.000){2}{\rule{4.577pt}{1.000pt}}
\put(726,804.42){\rule{9.154pt}{1.000pt}}
\multiput(726.00,803.92)(19.000,1.000){2}{\rule{4.577pt}{1.000pt}}
\put(764,804.42){\rule{9.154pt}{1.000pt}}
\multiput(764.00,804.92)(19.000,-1.000){2}{\rule{4.577pt}{1.000pt}}
\put(802,802.92){\rule{9.154pt}{1.000pt}}
\multiput(802.00,803.92)(19.000,-2.000){2}{\rule{4.577pt}{1.000pt}}
\put(840,799.92){\rule{8.913pt}{1.000pt}}
\multiput(840.00,801.92)(18.500,-4.000){2}{\rule{4.457pt}{1.000pt}}
\put(877,795.92){\rule{8.913pt}{1.000pt}}
\multiput(877.00,797.92)(18.500,-4.000){2}{\rule{4.457pt}{1.000pt}}
\multiput(914.00,793.69)(3.455,-0.462){4}{\rule{6.417pt}{0.111pt}}
\multiput(914.00,793.92)(23.682,-6.000){2}{\rule{3.208pt}{1.000pt}}
\multiput(951.00,787.69)(2.641,-0.475){6}{\rule{5.250pt}{0.114pt}}
\multiput(951.00,787.92)(24.103,-7.000){2}{\rule{2.625pt}{1.000pt}}
\multiput(986.00,780.68)(2.257,-0.481){8}{\rule{4.625pt}{0.116pt}}
\multiput(986.00,780.92)(25.401,-8.000){2}{\rule{2.312pt}{1.000pt}}
\multiput(1021.00,772.68)(1.860,-0.485){10}{\rule{3.917pt}{0.117pt}}
\multiput(1021.00,772.92)(24.871,-9.000){2}{\rule{1.958pt}{1.000pt}}
\multiput(1054.00,763.68)(1.860,-0.485){10}{\rule{3.917pt}{0.117pt}}
\multiput(1054.00,763.92)(24.871,-9.000){2}{\rule{1.958pt}{1.000pt}}
\multiput(1087.00,754.68)(1.554,-0.487){12}{\rule{3.350pt}{0.117pt}}
\multiput(1087.00,754.92)(24.047,-10.000){2}{\rule{1.675pt}{1.000pt}}
\multiput(1118.00,744.68)(1.501,-0.487){12}{\rule{3.250pt}{0.117pt}}
\multiput(1118.00,744.92)(23.254,-10.000){2}{\rule{1.625pt}{1.000pt}}
\multiput(1148.00,734.68)(1.195,-0.491){16}{\rule{2.667pt}{0.118pt}}
\multiput(1148.00,734.92)(23.465,-12.000){2}{\rule{1.333pt}{1.000pt}}
\multiput(1177.00,722.68)(1.261,-0.489){14}{\rule{2.795pt}{0.118pt}}
\multiput(1177.00,722.92)(22.198,-11.000){2}{\rule{1.398pt}{1.000pt}}
\multiput(1205.00,711.68)(1.065,-0.491){16}{\rule{2.417pt}{0.118pt}}
\multiput(1205.00,711.92)(20.984,-12.000){2}{\rule{1.208pt}{1.000pt}}
\multiput(1231.00,699.68)(0.977,-0.491){16}{\rule{2.250pt}{0.118pt}}
\multiput(1231.00,699.92)(19.330,-12.000){2}{\rule{1.125pt}{1.000pt}}
\multiput(1255.00,687.68)(0.934,-0.491){16}{\rule{2.167pt}{0.118pt}}
\multiput(1255.00,687.92)(18.503,-12.000){2}{\rule{1.083pt}{1.000pt}}
\multiput(1278.00,675.68)(0.820,-0.492){18}{\rule{1.942pt}{0.118pt}}
\multiput(1278.00,675.92)(17.969,-13.000){2}{\rule{0.971pt}{1.000pt}}
\multiput(1300.00,662.68)(0.740,-0.492){18}{\rule{1.788pt}{0.118pt}}
\multiput(1300.00,662.92)(16.288,-13.000){2}{\rule{0.894pt}{1.000pt}}
\multiput(1320.00,649.68)(0.716,-0.491){16}{\rule{1.750pt}{0.118pt}}
\multiput(1320.00,649.92)(14.368,-12.000){2}{\rule{0.875pt}{1.000pt}}
\multiput(1338.00,637.68)(0.620,-0.492){18}{\rule{1.558pt}{0.118pt}}
\multiput(1338.00,637.92)(13.767,-13.000){2}{\rule{0.779pt}{1.000pt}}
\multiput(1355.00,624.68)(0.540,-0.492){18}{\rule{1.404pt}{0.118pt}}
\multiput(1355.00,624.92)(12.086,-13.000){2}{\rule{0.702pt}{1.000pt}}
\multiput(1370.00,611.68)(0.500,-0.492){18}{\rule{1.327pt}{0.118pt}}
\multiput(1370.00,611.92)(11.246,-13.000){2}{\rule{0.663pt}{1.000pt}}
\multiput(1385.83,595.12)(0.491,-0.541){16}{\rule{0.118pt}{1.417pt}}
\multiput(1381.92,598.06)(12.000,-11.060){2}{\rule{1.000pt}{0.708pt}}
\multiput(1397.83,581.06)(0.489,-0.543){14}{\rule{0.118pt}{1.432pt}}
\multiput(1393.92,584.03)(11.000,-10.028){2}{\rule{1.000pt}{0.716pt}}
\multiput(1408.83,566.22)(0.481,-0.745){8}{\rule{0.116pt}{1.875pt}}
\multiput(1404.92,570.11)(8.000,-9.108){2}{\rule{1.000pt}{0.938pt}}
\multiput(1416.83,553.22)(0.481,-0.745){8}{\rule{0.116pt}{1.875pt}}
\multiput(1412.92,557.11)(8.000,-9.108){2}{\rule{1.000pt}{0.938pt}}
\multiput(1424.86,536.17)(0.424,-1.169){2}{\rule{0.102pt}{2.850pt}}
\multiput(1420.92,542.08)(5.000,-7.085){2}{\rule{1.000pt}{1.425pt}}
\put(1427.92,521){\rule{1.000pt}{3.373pt}}
\multiput(1425.92,528.00)(4.000,-7.000){2}{\rule{1.000pt}{1.686pt}}
\put(1431.42,508){\rule{1.000pt}{3.132pt}}
\multiput(1429.92,514.50)(3.000,-6.500){2}{\rule{1.000pt}{1.566pt}}
\put(1433.42,495){\rule{1.000pt}{3.132pt}}
\multiput(1432.92,501.50)(1.000,-6.500){2}{\rule{1.000pt}{1.566pt}}
\end{picture}



\vspace{1cm}
\caption{Intersections of periodic solutions of Henon--Heiles' system
with plane $X = 0,  \frac{dX}{dt} = \frac{dX}{dt}(X=0,Y,Y',H)$ at the
full energy ${1 \over 12}$.}
\end{figure}


   It is necessary to remark that for the value $H = {1 \over 6}$ we will
have a chaotic regime of a behavior of the system (Henon, Heiles, 1964).
So ${1 \over 24} \mbox{ and } {1 \over 12}$ are not physically small 
values.

The creation of the normal form and the normalizing transformation took
about 30 secs and brought 110 terms to the form and 1250 terms to the
transformation.



\section{The program }

   The earlier attempts of the author to realize the creation of
high enough orders of the normal form on algebraic level of
REDUCE-2 system were not very successful as a sequence of
small capacity of computers which the author used. And for 
realization of the described frame  
the NORT program (Edneral, Khrustalev, 1985, 1992), was created. The 
program NORT is written in Standard LISP and contains about  
1500 LISP operators.  It  is a package of procedures to treat 
truncated multivariate power series of arbitrary dimensions. In 
addition to the procedures for arithmetic operations there are 
special functions for the creation of normal forms and functions 
for substitutions, for calculation of roots (when it is possible), 
for differentiating, for printing and for inverting of multivariate 
power series.
    Coefficients of the truncated power series may be treated  
in three different (precision rational, floating point or an 
approximate rational (Rodionov, in progress) ) arithmetics 
according to the user's choice.
Separately user may choose a representation of numbers (in a 
floating point or in a rational form) for output. This output has 
a REDUCE  acceptable form,  but  REDUCE's  internal procedures are 
not called by the program. The program uses a recurrent polynomial 
representation separately in each order of the stored series. There is
a just compact use of memory. Unfortunately at this
moment the system has no friendly user interface and has a research
character. More that for nonlinear systems of ODE there is a so
wide variety that we have no clear ideas about proper realization of
the interface as it was possible for second orders systems 
(pay your attention to the nice package ALKAHEST III 
(Fitch, Norman, Moore, 1986)).

   The time of a garbage collection for above examples of second
order equations was smaller that 3\% of evaluation time. This
can characterize the NORT system as a system with a good enough 
internal LISP organization. 

   Regarding a speed of calculations under different arithmetics.
   The evaluation ov van der Pol's equation till 28 order
took under the rational arithmetic 21 minutes, under approximate
rational arithmetic it took 13 minutes 
with relative precision
about $10^{-7}$ and under float point arithmetic the evaluation 
took 161 secs.

\section{Conclusion}
  Here we can conclude that the creation of the high order normal 
forms allows to produce finite formulaes for a quantitative 
approximation of periodic solutions of nonlinear ODEs.


\paragraph{Acknowledgments:}
I am very grateful to A.D.Bruno for important advices  and useful 
discussions,  to O.A.Khrustalev for a concentration of 
my attention on the normal form method at the most beginning, to 
J.H.Davenport and J.P.Fitch for support at the stage of further
investigations. I would like also to express my acknowledgments
to H.Melenk who kindly introduced for my needs the "zeropoint" 
facility into the Groebner/REDUCE package and to S.Walcher who 
provided me with his newest papers.


\begin{thebibliography}{99}

\bibitem{AndersenGeer1983}
Andersen, G.M., Geer, J.F. (1983). Power series expansions for
the frequency and period of the limit cycle of the van der Pol
equation. {\em SIAM J. Appl. Math.\/} {\bf 42}, 678--693.

\bibitem{ArnoldAnosov1988}
Arnold, V.I., Anosov, D.V. (Eds.) (1988). Dynamical Systems, I. 
{\em Encyclopaedia of Mathematical Sciences\/}. Springer--Verlag,
N.Y.

\bibitem{BoegeGebauerKredel1986}
Boege, W., Gebauer, R., Kredel, H. (1986). Some
Examples for Solving Systems of Algebraic Equations by Calculating
Groebner Bases. {\em  J. Symbolic Computation\/} {\bf 1}, 83--98.

\bibitem{Bruno1971}
Bruno(Brjuno), A.D. (1971)
{\em Analytical form of differential equations. I.\/} Trans. Mosc. 
Mat. Soc.,  {\bf 25}, 131--288.

\bibitem{Bruno1972}
Bruno(Brjuno), A.D. (1972)
{\em Analytical form of differential equations. II.\/} Trans. Mosc. 
Mat. Soc. {\bf 26}, 199--239.


\bibitem{Bruno1989}
Bruno, A.D. (1989)
{\em Local Method in Nonlinear Differential Equations. 
Part  I - The Local Method of Nonlinear Analyses of  Differential 
Equations,  Part  II - The Sets of Analyticity of  a  Normalizing 
Transformation\/}. Springer Series in Soviet Mathematics.  ISBN 
3-540-18926-2,  pp. 370.

\bibitem{Bruno1990}
Bruno, A.D. (1990)
{\em The restricted three-body problem: Plane periodic orbits\/}. 
Moscow, Nauka. ISBN 5-02-000683-1, 295. (in Russian).

\bibitem{Bruno1993} 
Bruno, A.D. (1993) Bifurcation of the Periodic
Solutions in the Case of a Multiple Pair of Imaginary Eigenvalues.
{\em Selecta Mathematica {\em formerly} Sovietica\/} {\bf 12}, 
\# 1, 1--12.

\bibitem{BrunoWalcher} 
Bruno, A.D., Walcher, S. (to be printed) Symmetries and convergence of
normalizing transformations.
{\em J. Math. Anal. Appl.\/}

\bibitem{Buchberger1985}
Buchberger, B. (1985).
Groebner Bases:  An Algorithmic Method in Polynomial Ideal Theory.
{\em Progress, directions and open problems in multidimensional
systems theory\/}. Dordrecht: Reidel, edited by N.K.Bose, 184--232.

\bibitem{Deprit1969}
Deprit, A. (1969) Canonical transformation depending on a small
parameter. {\em Celestial Mechanics} {\bf 1}, \# 1, 12--30.

\bibitem{GodziewskiMaciejewski1990}
Godziewski, K., Maciejewski,A.J. (1990)
{\em Celest. Mech. Dyn. Astron.\/} {\bf 49}, 1.

\bibitem{Ito1989}
Ito, H. (1989) Convergence of Birkhoff normal forms for integrable systems.
{\em Comment. Math. Helv.\/} {\bf 64}, 412--461.

\bibitem{Ito1992}
Ito, H. (1992) Integrability of Hamiltonian systems and Birkhoff normal 
forms in the simple resonance case.
{\em Mth. Ann.\/} {\bf 292}, 411--444.

\bibitem{EdneralKhrustalev1985} 
Edneral, V.F., Khrustalev, O.A. (1985).
The Normalizing Transformation for Nonlinear Systems of ODE.  The 
Realization  of  the  Algorithm. {\em Proceedings of International 
Conference on Computer Algebra and its Application in  Theoretical  
Physics (USSR, Dubna, September 1985)\/}. Dubna: JINR publ.,  
219--224. In Russian.

\bibitem{EdneralKhrustalev1992}
Edneral, V.F., Khrustalev, O.A. (1992). Program  for  Recasting  ODE 
Systems in the Normal Form. {\em Sov. J. Programmirovanie\/}, \# 5, 
73--80. In Russian.

\bibitem{Edneral1993}
Edneral, V.F. (1993). Computer Generation of Normalizing Transformation  
for  Systems  of  Nonlinear   ODE. 
{\em  Proceedings  of the 1993 International Symposium on Symbolic  
and Algebraic Computation (Kiev,  Ukraine, July 1993)\/}. New York: ACM 
Press, edited by M.Bronstein, 14--19.
 
\bibitem{FitchNormanMoore1986}
Fitch, J.P., Norman, A.C., Moore, P.M.A. (1986). ALKAHEST III: Automatic 
analysis of periodic weakly nonlinear ODEs. {\em  Proceedings of
SYMSAC 86.\/} Waterloo, 34---38.

\bibitem{GuckenheimerHolmes1986}
Guckenheimer, J., Holmes, P. (1986) {\em Nonlinear Oscillations, 
Dynamical Systems and Bifurcations of Vector Fields}.  Springer
--Verlag, N.Y.

\bibitem{HassardKazarinoffWan1981}
Hassard, B.D., Kazarinoff, N.D., Wan, Y.-H. (1981)
{\em Theory and Applications of Hopf Bifurcation.\/} Cambridge Univ. 
Press. 280.

\bibitem{Hearn1987}
Hearn, A.C. (1987). REDUCE. User's manual. {\em  Rand Publication\/} 
CP78.

\bibitem{HenonHeiles1964}
Henon, M., Heiles,C. (1964). The Applicability 
of  the Third Integral Of Motion: Some  Numerical  Experiments. 
{\em Astronomical J.\/} {\bf 69}, 73--79.

\bibitem{Hori1966}
Hori, G.I. (1966).Theory of general perturbations with unspecified
canonical variables. {\em J. Japan Astron. Soc.\/} {\bf 18}, \# 4, 
287--296.

\bibitem{Mersman1970}
Mersman, W.A. (1970)  A new algorithm for Lie transformation. 
{\em Celestial Mechanics} {\bf 3}, \# 1, 81--89.

\bibitem{RandArmbruster1987}
Rand, R., Armbruster D. (1987) {\em Perturbation methods, Bifurcation 
Theory and Computer Algebra\/}. Springer--Verlag, N.Y.

\bibitem{Rodionov}
Rodionov, A.Ya. (in progress). A Program for Rational Arithmetic 
RATAR. (In Russian).

\bibitem{ShevchenkoSokolsky1993}
Shevchenko, I.I., Sokolsky, A.G. (1993) Algorithms for normalization of
Hamiltonian systems by means of computer algebra. {\em Comp. Phys. Comm.}
{\bf 77}, 11---18.

\bibitem{Vallier1993}
Vallier, L. (1993) An Algorithm for the Computation of Normal Forms and 
Invariant Manifolds.
{\em  Proceedings  of the 1993 International Symposium on Symbolic  
and Algebraic Computation (Kiev,  Ukraine, July 1993)\/}. New York: ACM 
Press, edited by M.Bronstein, 225---233.

\bibitem{Walcher1991}
Walcher, S. (1991) On differential equations in normal form. 
{\em Math. Ann.\/} {\bf 291}, 293---314.

\bibitem{Walcher1993}
Walcher, S. (1993) On Transformations into Normal Form. {\em J. Math. 
Analysis and Appl.\/} {\bf 180}, 617---632.

\end{thebibliography}

\end{document}
